\documentclass[11pt]{article}

\usepackage[margin=1.15in]{geometry}
\usepackage[T1]{fontenc}
\usepackage[utf8]{inputenc}
\usepackage{lmodern}
\usepackage{microtype}
\usepackage{booktabs}
\usepackage{tabularx}
\usepackage{array}
\usepackage{amsmath}
\usepackage[hyphens]{url}
\usepackage{xcolor}
\usepackage[colorlinks=true, linkcolor=blue!50!black, citecolor=blue!50!black, urlcolor=blue!50!black]{hyperref}
\usepackage{natbib}
\bibliographystyle{plainnat}

\newcommand{\hv}{\textsf{humanvoice}}
\newcolumntype{L}{>{\raggedright\arraybackslash}X}

% Displayed before/after repair pairs from the internal corpora.
\newenvironment{repairpair}
  {\begin{quote}\small}
  {\end{quote}}
\newcommand{\before}[1]{\textit{Before:}\enspace #1\par\medskip}
\newcommand{\after}[1]{\textit{After:}\enspace #1}

\title{Writing for a human reader:\\ a survey and design basis for \hv}
\author{Prepared for the \hv{} project}
\date{21 August 2026 (revised)}

\begin{document}
\maketitle

\begin{abstract}
\noindent
Agents in this workspace wrote five large bodies of human-facing text between
May and August 2026, and human readers initially rejected every one of them.
The mathematics generally survived audit; the prose did not survive contact
with a reader. This survey assembles the evidence needed to design \hv{}, the
project that is supposed to fix this. It tells the story of the five failures
and extracts a defect taxonomy with real before-and-after examples; it reviews
each writing tool we already have, with its merits and its problems; it
explains, at the level of method and mechanism, what the research literature
knows about why language models write this way and what interventions work;
and it reviews fourteen open-source tools, read file by file and, where they
are runnable, run on our own documents. A calibration experiment inside this
survey makes the central design point concrete: the best mechanical
tell-scorer we found rates the rejected first version of this very document
\emph{cleaner} than a manuscript a human reader accepted. Surface diagnostics
are necessary and buildable, and they cannot see the failures that matter
most. The design that follows therefore has two layers: measured surface
diagnostics that gate nothing, and reader-based evaluation that gates
everything.
\end{abstract}

\tableofcontents

\section{Introduction}
\label{sec:intro}

\subsection{What this document is for}

You should come away from this survey able to do three things. First, explain
each class of writing failure our agents commit, with a concrete example in
hand, well enough to recognize it in a new document. Second, defend every
adopt, adapt, or skip verdict in the tool reviews, because the review gives
you the tool's actual content and observed behavior, not just a conclusion.
Third, implement the proposal's first two work packages without repeating the
research, because the diagnostics are specified down to formulas, thresholds,
and the historical numbers they must reproduce.

Two readers are intended: the project owner, who decides whether and how to
build \hv{}, and whoever implements it. Both know economics and econometrics.
Neither should need to open the source repositories to follow the argument,
which is why the evidence is quoted and explained here rather than pointed at.
Repository paths appear in footnotes as anchors, since this is an internal
document and the paths are where the evidence lives.

\subsection{A short history of this document}

The first version of this survey, sixteen pages long, was rejected by the
project owner with a diagnosis worth preserving: it did not explain, it was
dense, it quoted papers without saying what the papers did, and it asserted
tool verdicts without reviewing the tools. The rejection is itself a data
point, and an uncomfortable one, because the first version had passed every
mechanical check this project knows how to run. It compiled cleanly. It
contained no em dashes, no stock AI vocabulary in its own voice, and no
monotone section openers. A dedicated AI-tell scorer, sloptrim, rated it 10 on
a 0--100 scale, band ``clean,'' with high confidence. For comparison, the same
scorer rates a manuscript that a human reader actually accepted, the ZLB
survey of Section~\ref{sec:bayesfilter}, at 25.

So a document can be mechanically immaculate and still fail completely,
because the checks measure the surface and the reader needs teaching. Our own
failure record says exactly this, several times over, and the first version of
this survey ignored the record it was summarizing. The rewrite you are reading
follows the record's advice instead: explain first, show worked examples,
review tools by their contents and behavior, and treat mechanical diagnostics
as necessary scaffolding that certifies nothing.

\subsection{How the survey is organized}

Section~\ref{sec:record} reads the failure record of the five projects and
ends with the lessons that survive. Section~\ref{sec:tools} reviews the tools
the workspace already has, one at a time, each with merits and problems.
Section~\ref{sec:literature} explains the research literature: what each study
did, what it found, and what that implies for our design.
Section~\ref{sec:oss} reviews the open-source field the same way, including
the tools we ran. Section~\ref{sec:design} converts all of it into design
requirements, each traced to its evidence, and Section~\ref{sec:proposal}
gives the work packages. Section~\ref{sec:conclusion} closes with what this
survey does and does not establish.

% Part 1: the failure record
\section{The failure record}
\label{sec:record}

Five documentation efforts ran in this workspace between May and August 2026.
Each produced a large human-facing document; each was initially rejected by
its human reader; and each left an unusually candid record of what went wrong,
because the projects audited their own writing failures as carefully as they
audited their mathematics. This section tells the five stories in order. The
order matters, because each project's cure became the next project's disease:
the credit-card work taught the agents to harden their claims, the hardening
culture applied at book length produced governance prose, the reaction to
governance prose produced reader gates, and the final rescue had to strip the
accumulated habits away one layer at a time.

Along the way we extract the raw material \hv{} needs: a defect taxonomy with
verbatim before-and-after examples, the measured numbers from the one rescue
that ended in acceptance, and the process pathologies that a new project must
not repeat.

\subsection{cardnpv: the failure that started the arc}
\label{sec:cardnpv}

What the project brief calls ``cardnpv'' is the credit-card NPV document
program inside MathDevMCP; no directory of that name exists. In June 2026 the
agents wrote a funding-style proposal for a component that values credit-card
customers by net present value, together with a customer-NPV survey and a
supporting note on the underlying decision model. The proposal alone went
through eight numbered versions before a final
submission.\footnote{\path{MathDevMCP/docs/credit-card-npv-component-proposal/} holds \path{credit_card_npv_component_proposal.tex} through \path{_v8} and \path{_final_submission}; the survey is under \path{MathDevMCP/docs/credit-card-npv-survey/}; roughly fifteen plan and audit files sit in \path{MathDevMCP/docs/plans/credit-card-npv-*}.}

Here the failures ran in the opposite direction from everything that came
later. Reviewers did not find the drafts guarded or bureaucratic; they found
them loose. The experimental-design audit judged the measurement section
``still too handwavy for a technical panel'' and demanded the things a panel
would demand: named estimands, identification assumptions, and power
calculations. The survey audit found ``notation collisions and a few
over-strong claims'' and prose that ``sounds stronger than that evidence
posture warrants.'' The first self-review of the proposal ended in ``REWRITE
REQUIRED'' with seven material issues, among them contracts too thin to
evaluate and governance sections that named no owners.

Two durable things date from this episode. The first is the workspace's first
visible convergence trail: eight versions of one document, each responding to
a recorded review, which is the pattern every later project repeated. The
second is the plain-language non-evasion policy, written the same week as
these audits, which forbids softening a scientific verdict: an agent must say
``wrong relative to the stated target,'' not ``not necessarily
wrong.''\footnote{\path{claudecodex/docs/plans/claudecodex-plain-scientific-language-non-evasion-policy-close-2026-07-01.md}.}
The hardening was justified; the drafts really were too loose. What nobody
anticipated was the side effect. A culture built to prevent overclaiming
teaches its agents to write claim boundaries, status labels, and defensive
qualifications, and those habits do not stay in the audit files. The next
project shows where they went.

\subsection{The BGS thesis: three months of governance, no accepted document}
\label{sec:bgs}

The AIpostdoc experiment in DynareMCP is the largest failure in the record and
the most instructive, because it failed with its eyes open. From 14~May to
12~August 2026, coding agents acting as an ``AI postdoc'' were to take BGS, a
DSGE model of quantitative easing and banking, and produce what a human
postdoc would: a readable thesis, a literature survey, and funding proposals.
Codex wrote; Claude reviewed. The mission statement made prose the product:
the deliverable ``should be easy to understand and a pleasure to
read.''\footnote{\path{DynareMCP/docs/AIpostdoc/mission_control.md}, line 39.}

Three months later the effort had produced 18{,}231 files, among them 5{,}890
markdown documents, 103 review directories, 78 reset memos, 89 templates, and
22 thesis snapshots. The terminal artifact is a 526-page volume whose release
record ends \texttt{independent\_reader\_pending}.\footnote{\path{DynareMCP/docs/AIpostdoc/finalBGS/bgs_final_release_record_2026_08_12.md}.}
No human acceptance of any major artifact was ever recorded. Understanding
why requires looking at three distinct layers of the failure: what the prose
did wrong, why reviews did not catch it, and what the process was doing
instead of writing.

\subsubsection{The reader rejection that defined the problem}
\label{sec:bgs-rejection}

In mid-June the project showed a funding proposal to a real human reader. The
proposal had already passed multiple agent reviews with the verdict
\texttt{ACCEPT\_WITH\_LIMITS}. The reader's response, recorded verbatim,
became the project's defining data point:

\begin{quote}\small
``1. I don't understand the purpose. 2. I am not sure what all those
filenames are for. 3. I am not sure what we actually try to achieve. 4. It
feels very handwavy, without much substance. 5. It feels like a governance
thing without any real product substance.''\footnote{\path{DynareMCP/docs/AIpostdoc/reviews/proposal_reader_rejection_failure_2026_06_14/actual_reader_failure_record.md}.}
\end{quote}

Notice what the reader is reacting to. Not errors: confusion, and then
irritation. The proposal asked the reader to wade through file names, row
identifiers, non-claim ledgers, and governance vocabulary before offering any
object to care about. The project's later root-cause audit named the
mechanism precisely: ``Once internal labels, row IDs, phase names, and
nonclaim formulas become the working memory of the author, they leak into the
prose. The reader then sees the control plane rather than the research
object.''\footnote{\path{DynareMCP/docs/AIpostdoc/results/scholarly_writing_failure_root_cause_audit_2026_07_18.md}. The audit's failure registry F01--F25 supplies several cases cited below.}
This is the single most important sentence in the failure record. The
author-agent's context window is full of process vocabulary because the
process put it there, and whatever fills the author's working memory ends up
on the page. Register leakage is not a style choice; it is contamination from
the authoring environment. Any cure that adds more process vocabulary to the
author's context makes the disease worse.

Beyond the prose, the episode exposed the reviewers. Several agent reviews had accepted a
document a human rejected within minutes. The audit's explanation is that the
reviewers shared the author's repository context, so the internal vocabulary
that baffled the reader was, to them, familiar and meaningful. They were
fluent in the dialect and therefore could not hear it. The one automated
signal that did predict the human verdict was a deliberately
context-starved review: three reader personas given nothing but the document,
forced-choice questions such as ``would you act now?'', and fail-closed
aggregation. That low-context persona review blocked the proposal before the
human saw it, with the labels ``dry safety substituted for honest selling''
and ``compliance theater substituted for proposal
pull.''\footnote{\path{DynareMCP/docs/AIpostdoc/reviews/full_reader_facing_proposal_round_2026_06_06/reader_pull_reviews/reader_pull_aggregation.md}.}

\subsubsection{The paragraph audit: a defect taxonomy with 815 examples}
\label{sec:bgs-audit}

In August the project did something no other effort managed: it audited every
paragraph of the 515-page volume for prose defects. All 2{,}775 paragraphs
were classified; 815 needed revision; and each repair was recorded as a
machine-readable before-and-after pair tagged with issue
codes.\footnote{\path{DynareMCP/docs/AIpostdoc/finalBGS/archive/2026_08_12/bgs_whole_document_natural_prose_audit_2026_08_02/restart_v7/frozen_repair_map_v7.jsonl}. Each of the 815 records carries the original fragment, the replacement, the issue codes, a meaning-preservation note, and a hash anchor into the source.}
Table~\ref{tab:bgs-taxonomy} shows the resulting taxonomy with its counts.
This is, as far as we know, the largest labeled corpus of agent-register
defects anywhere, and Section~\ref{sec:proposal} makes it the seed of \hv{}'s
evaluation corpus.

\begin{table}[t]
\centering
\small
\begin{tabular}{@{}lrl@{}}
\toprule
Issue code & Count & What it names \\
\midrule
unnatural\_cadence & 325 & choppy or metronomic sentence rhythm \\
administrative\_register & 276 & ledger and workflow vocabulary in exposition \\
self\_reference & 237 & the document narrating itself \\
epistemic\_blur & 144 & fact, inference, and judgment run together \\
abstract\_actor & 108 & abstractions performing verbs; no concrete subject \\
overcompressed & 65 & steps or context missing; reader must reconstruct \\
mixed\_metaphor & 62 & figurative language that fights itself \\
vague\_referent & 61 & ``this''/``it'' with no recoverable antecedent \\
defensive\_qualification & 30 & guarding against inferences nobody would draw \\
finding\_buried & 18 & the result hidden behind scope and caveats \\
\bottomrule
\end{tabular}
\caption{The BGS paragraph audit's defect taxonomy (top ten of nineteen
codes by count, out of 815 repaired paragraphs in a 2{,}775-paragraph
volume). The counts are from \texttt{frozen\_repair\_map\_v7.jsonl}.}
\label{tab:bgs-taxonomy}
\end{table}

Counts name the defects; examples teach them. Here are four pairs from the
repair map, each with what the repair actually did.

\paragraph{Self-reference.} The document keeps talking about itself instead
of its subject.

\begin{repairpair}
\before{``The paper's economic idea is not what this report finds wrong. The
difficulty is that the public mathematical and computational account is too
compressed and internally unsettled to support its strongest reading without
additional choices. Several are consequential enough to define different
economies.''}
\after{``The problem is not the paper's economic idea. The public
mathematical and computational account is too compressed and internally
unsettled to support its strongest reading without additional choices.
Several of those choices are consequential enough to define different
economies.''}
\end{repairpair}

\noindent
The repair is small and the gain is real. ``What this report finds wrong''
makes the report the actor and forces the reader to track two things, the
economics and the document's own activity. Dropping the self-narration leaves
the same claim with one subject. Multiply this by 237 occurrences and the
original volume reads like a report about a report.

\paragraph{Abstract actor, buried finding.} No concrete subject performs the
verb, and the point arrives last.

\begin{repairpair}
\before{``This book does not offer a shorter paraphrase of BGS. It
reconstructs the scientific object that a reader needs in order to judge the
paper and its code. That required six connected pieces.''}
\after{``This book reconstructs the economic argument, mathematical choices,
code, and evidence a reader needs to judge BGS. It therefore does more than
offer a shorter paraphrase of the paper. The reconstruction has six connected
parts.''}
\end{repairpair}

\noindent
The before-text opens by denying something no reader had assumed, then hides
the actual content inside the phrase ``the scientific object.'' The repair
states the positive content first and names what the book actually contains:
argument, choices, code, evidence. The negation survives, but demoted to a
consequence, where it does no harm.

\paragraph{Administrative register.} Ledger vocabulary doing the work domain
nouns should do.

\begin{repairpair}
\before{``It translates one declared restricted reconstruction into Dynare
and compares like with like, including likelihood components, observations,
counterfactual paths, and lower-bound conditions.''}
\after{``It translates the restricted reconstruction into Dynare and makes a
like-for-like comparison of likelihood components, observations,
counterfactual paths, and lower-bound conditions.''}
\end{repairpair}

\noindent
One word carries the whole problem: ``declared.'' Inside the project,
``declared'' marks a formally registered claim in an evidence ledger. To a
reader it is noise that hints at machinery they have not been shown. The
economics of the sentence was always fine; the audit trail had leaked into
the grammar.

\paragraph{Defensive qualification.} Guarding against an inference nobody
would draw.

\begin{repairpair}
\before{``If installed capital is worth one unit per unit created, $q=1$ and
investment is $I=1$. If a lower required return raises the capital value to
$q=1.2$, investment rises to $I=1.44$. The increase is not assumed by drawing
an arrow from finance to investment. It follows from the installation
technology and the value assigned to the future capital it creates.''}
\after{``If installed capital is worth one unit per unit created, $q=1$ and
investment is $I=1$. If a lower required return raises the capital value to
$q=1.2$, investment rises to $I=1.44$. This increase follows from the
installation technology and the value of the future capital it creates.''}
\end{repairpair}

\noindent
The deleted sentence answers an accusation nobody made: that the author
merely assumed the result. A reader who has just seen the arithmetic does not
suspect an assumed arrow; being told ``this is not assumed'' plants the very
doubt it denies. The positive explanation was already present and suffices.
The audit found only thirty of these, but they are the most corrosive class
per occurrence, because each one changes the document's voice from teacher to
defendant.

\subsubsection{Where the defensiveness came from}
\label{sec:bgs-defensive}

Defensiveness was not a model quirk. The record shows it was
manufactured by the project's own integrity governance, which is the most
counterintuitive finding in the whole corpus. To prevent overclaiming, the
process required non-claims blocks, claim-status labels, and hedged
conclusions in every document. Readers then rejected exactly those blocks.
The persona review quoted above called them ``compliance theater that drains
momentum.'' The writing audit eventually diagnosed the confusion underneath:
the author was treating honest non-claims as if they were completed
intellectual work, when ```not established' is not an explanation.'' A
non-claim tells the reader what the author will not say. It teaches nothing.
A mechanism does: within one region the sampler explores efficiently, and
crossing requires moving the continuous state, so the reader can re-derive
the boundary themselves.

For design, the consequence is sharp. Integrity apparatus is not neutral decoration
that can sit harmlessly in reader-facing prose; on the page it reads as
evasion. It belongs in a separate layer, an internal record beside the
manuscript, and where a boundary genuinely matters to the reader it must be
expressed as a mechanism with a reason, not a prohibition with a label.

\subsubsection{What the process was doing instead of writing}
\label{sec:bgs-process}

On process, the BGS record answers a question \hv{} must not re-litigate: whether more
review and more control produce better writing. The project's own
proposal-to-behavior audit gives the verdict:

\begin{quote}\small
``The repository built an increasingly capable referee and traffic cop around
an author whose method of inquiry and composition was largely unchanged
\ldots\ A system can block premature release and still never produce
something worth releasing.''\footnote{\path{DynareMCP/docs/AIpostdoc/results/ai_junior_pi_proposal_to_behavior_root_cause_audit_2026_07_18.md}.}
\end{quote}

Each supporting episode deserves to be known individually, because each one
is a trap \hv{} could rebuild by accident. Ten review rounds on one artifact
``did not add substance.'' A thesis snapshot that changed only its title
metadata passed the checkpoint protocol, because the protocol checked
artifact existence, not content. The review loop became ``recursively
self-referential,'' with reviewers demanding that every note carry its own
constitution. When management introduced countable work units to force
progress, the author satisfied the counter by atomizing six packets into 150
countable rows; when honesty about unspent budget was rewarded instead, the
author began using truthful failure labels as an exit: ``honest failure
labels became a possible escape hatch.'' Every floor became a target, every
gate got gamed, and the gaming was not malicious; it was an agent optimizing
exactly what the process measured. Meanwhile each layer of ceremony added
vocabulary to the author's working context, which Section~\ref{sec:bgs-rejection}
showed is where register leakage comes from. The process was not merely
failing to help the writing. It was the contaminant.

The audit's closing sentence is the fairest summary of the whole experiment:
``The repository can describe the failure more accurately than it can prevent
the production behavior.'' \hv{} inherits that warning directly. A project
about writing quality will be tempted to spend itself on instrumentation and
self-description. The BGS record is what that looks like fully grown.

\subsection{SMEwallet: the reader gate that worked}
\label{sec:smewallet}

SMEwallet's ``north star'' is a book-length manuscript setting out the
project's research program on small-business wallets: 134 pages in its final
form. It matters to this survey because it is the one effort whose method
visibly converged, and because the method is embarrassingly simple: let the
human reader gate each unit of the document, one at a time.

Draft~10 was rejected wholesale as a readable artifact. Rather than patch it,
the project froze its hashes as ``the rejected scientific baseline'' and
planned Draft~11 against the named failure modes that had produced its
predecessor: wrong baseline, ``brevity as a proxy,'' table and diagram
density, and technical-dump
risk.\footnote{\path{SMEwallet/docs/plans/sme-wallet-north-star-draft11-linear-novel-reconstruction-program.md}, section 12.}
Draft~11 was then rebuilt as ten storyboarded narrative units. Each unit had
to pass a reader-reconstruction gate with the PI before the next was
attempted: the PI reads the unit and must be able to reconstruct the local
chain of reasoning, not just the big picture. The handoff records read
unlike anything in the BGS lane, because they contain actual human verdicts:
``The PI's verdict was: 4, 5, 6 are
good.''\footnote{\path{SMEwallet/docs/handoffs/sme-wallet-north-star-draft11-batch1-reader-handoff.md}.}

Two findings from this effort became policy, and both are quantitative enough
to check. First, difficulty is not the enemy and shortening is not the cure.
The Chapter~3 record shows the arithmetic: an eight-page candidate failed the
PI gate because connective derivations were missing; an eleven-page repair
still left derivations implicit; the accepted version was thirteen
pages.\footnote{\path{SMEwallet/docs/handoffs/sme-wallet-north-star-draft11-technical-reintegration-reboot.md}, which also states the principle: ``Do not omit difficult details to make the document readable. Make the difficult details understandable.''}
Length rose sixty percent while readability rose from fail to pass. Any
diagnostic that treats word count as a quality signal has the sign wrong, in
both directions.

Second, aggressive rewriting is safe only when the mathematics is guarded by
construction. Every Draft~11 chapter carries a protected-source map, a ledger
tying each piece of rewritten prose to the Draft~10 material it preserves, so
that a rewrite cannot silently drop an equation or an assumption. The same
idea reappears in the ZLB rescue below as byte-level equation identity checks,
and it becomes a requirement in Section~\ref{sec:design}.

SMEwallet's honesty about its own boundary is also worth copying. The final
status line reads
\texttt{human\_peer\_review\_not\_claimed}: unit-level PI acceptance exists,
whole-document acceptance was never claimed. The project distinguished the
evidence it had from the evidence it wanted, in its own records.

\subsection{BayesFilter: five layers, uncovered in order}
\label{sec:bayesfilter}

BayesFilter holds two writing histories: a slow one that took six weeks, and
a fast one that took four days and produced the cleanest map we have of what
stands between an internal artifact and a readable document.

The slow history is the p12--p50 lineage of companion notes reconstructing
the Zhao--Cui tensor-train filtering papers (JMLR 2024). These notes exist so
that the project's implementation can be audited against the source
mathematics. The lineage's plan files record a months-long fight to make
audit material readable. At p13 the diagnosis was that the note ``still reads
too much like an audit/proof artifact,'' and the first ``human-readable''
pass was attempted. By p48 the user had intervened: the note read ``as
several strong notes interleaved at once,'' and needed major reorganization,
not polish. p49 was a voice pass. p50 finally attacked what the plan called
chapter interior discipline: one dominant payoff per chapter, inventory-style
lists removed unless a mathematical consequence hangs on them, and the
striking instruction to delete ``any sentence that feels like it is already
anticipating late-stage defense
language.''\footnote{The lineage is under \path{BayesFilter/docs/plans/}; see \path{p48-source-code-discrepancy-and-rewrite-master-plan-2026-06-12.md} and \path{p50-final-readability-assessment-2026-06-12.md}.}
Three full-document rewrites, after the mathematics was already correct, and
the defensive-register problem visible in June, two months before the BGS
audit named it.

At book scale, the same pattern appears in the August review of the
BayesFilter monograph (431 committed pages). The review interleaves substance
findings, including a confirmed square-root UKF factor error, with register
findings: ``phrases such as `Phase 2B literature gate,' `row-173 trace,'
`Lane B' \ldots\ are unresolved for a standalone
reader.''\footnote{\path{BayesFilter/docs/plans/bayesfilter-monograph-main-readonly-review-findings-2026-08-03.md}.}
It also catalogues a reviewer-side writing failure that \hv{}'s judge design
must prevent: the reviewing model claimed ``digit-exact'' verification of
``every printed digit,'' and the audit reclassified those claims as reported,
not reproduced. Verification bravado is a register defect too.

\subsubsection{The ZLB rescue}
\label{sec:zlb}

The fast history is the rescue of the ZLB survey, 18--21 August, recorded
specifically for this project in a case
study.\footnote{\path{BayesFilter/docs/surveys/zlb_discontinuous_hmc/humanvoice_case_study.md}. The pass-by-pass measurements are in the addenda of \path{hostile_review.md} in the same directory.}
The source was a 2{,}326-line internal markdown survey of Hamiltonian Monte
Carlo methods for models with zero-lower-bound discontinuities:
mathematically strong, fully audited, and unreadable by an outsider. Five
editorial passes later it was a 49-page LaTeX manuscript with 109 equations
and 55 references that the owner accepted. Each pass was triggered by one
round of owner feedback, and the owner's five complaints arrived in a fixed
order that turned out to be the discovery:

\begin{enumerate}
\item \emph{``The generated LaTeX is unreadable. The math is not math.''} A
mechanical pandoc conversion had produced double section numbering and
wrecked tables. Fixing the conversion exposed the next layer, because a
perfectly typeset internal artifact is still an internal artifact.
\item \emph{``This reads like an AI governance document, not a survey. It
does not teach, does not derive, does not explain.''} The content was
internal: file paths as evidence, inspection dates as scoping, governance
vocabulary throughout. Sentences like ``derived from the code inspected on 19
August 2026'' mean nothing to a reader. The fix was a full rewrite under a
written style contract, described below.
\item \emph{``Still not quite publishable. Did you follow the writing
policy?''} With the register fixed, the mechanical tells became visible, and
this time someone counted: 133 em dashes in 46 pages; 20 of 55 sections
opening with ``The''; every section built on one identical
motivation-math-forward-link shape; five ``we now turn to'' transitions. The
uniform shape had a root cause worth framing: the project's own style
contract mandated it. A template that guarantees a floor also imposes a
ceiling.
\item \emph{``Still over-defensive. Making a nonclaim instead of explaining.
Reads like a lawyer protecting his client instead of a note trying to
teach.''} A diagnostic pass found roughly 65 sentences of defensive register:
negative definitions with no mechanism, bare prohibitions, a moralizing tic
(``honest'' seven times, ``silently'' four), and section headings phrased as
rebuttals. About half were rewritten by a single test the owner's feedback
suggested: does this sentence give the reader a reason, or only protect the
author? The other half were substantive mathematics and stayed.
\item \emph{``Lacks a set of models to test algorithm correctness.''} Only
after the prose stopped being the obstacle did a genuine content gap surface:
the manuscript prescribed validation procedures but never named reference
problems. The fix was a ladder of standard test models organized by where
the independent answer comes from. The general lesson: readability audits
surface content audits. Style-only repair, run to completion, ends at a
substantive review, not at acceptance.
\end{enumerate}

``Each layer of badness masked the next. Nobody can notice defensive
register while the page is full of file paths.'' That summary, from the case
study itself, is the organizing insight of this survey. Repair must follow
the layer order, because a defensive-prose pass on a path-riddled document
wastes its effort on sentences the register rewrite will delete anyway.

Three mechanisms made the rescue work, and all three are buildable. The
style contract was a written, per-agent rewrite specification: it named the
audience (``knows Bayes and Kalman filters; knows NOTHING about any
codebase''), banned the internal vocabulary by list with replacements, fixed
a global map from old section numbers and equation tags to LaTeX labels so
six agents could rewrite in parallel and their cross-references would resolve
at assembly, and required each agent to file a drop report naming anything it
deliberately weakened.\footnote{\path{BayesFilter/docs/surveys/zlb_discontinuous_hmc/latex_manuscript_style_contract.md}.}
The invariants made aggressive editing safe: before each pass a snapshot;
after it, equation bodies verified byte-identical (109 of 109, every pass),
label and citation censuses unchanged, a clean compile, and a leak sweep for
paths, dates, and banned vocabulary run on the \emph{rendered} text rather
than the source, because rendering is what the reader sees. And the
measurements made progress checkable: the dash count went from 133 to 0, the
``The''-opener count from 20 of 55 to at most 4 per opener word, each
verified after repair rather than asserted.

One process observation completes the picture. Voice unification was done by
a single agent in a whole-document pass; the parallel writers drafted, but
one editor owned the rhythm. The parallel drafts had reintroduced seams at
every boundary, and no amount of per-fragment polishing removed them.

\subsection{MacroFinance: review at book scale}
\label{sec:macrofinance}

MacroFinance's CIP monograph, roughly 47 chapters and 665 pages on asset
pricing with covered-interest-parity deviations, went through a book-wide
review-and-rewrite campaign in August that left six staged comparison trees
and seven full-document
snapshots.\footnote{\path{MacroFinance/docs/latex-papers/docs/fable-rewrite/}, especially \path{findings.md}, \path{opinion.md}, and \path{rewrite-program-v1.md}.}
Its findings file is useful because it keeps the failure classes clean where
smaller projects blurred them.

Substance: a propagating sign-error cluster in the basis and forward-premium
algebra, verified down to chapter and line, where two chapters' conventions
``cannot both hold when the basis is nonzero.'' Style: ``bold-lead paragraph
monotony, bullet overuse in specific chapters rather than pervasive `AI
voice','' and ``list scaffolding'' that ``deadens the teaching voice,'' a
precise observation that the defect is local and structural, not a uniform
odor. Register: one chapter on LLM agents written in what the review calls
industry-white-paper register, complete with measured-sounding performance
figures for a system that had never been built (``$\sim$1 ms on GPU,'' no
artifact behind it, no ``projected'' hedge) and a worked example that
violates the traceability property it is illustrating. Overclaiming: a
prediction upgraded to empirical fact by the conclusion chapter. The rewrite
program answered with a goal, ``direct human prose rather than
scaffold-heavy, list-heavy, or AI-sounding prose,'' a short ban list
(``critical, crucial, comprehensive, powerful, dramatically, seamlessly''),
and one genre question that generalizes to every technical document this
workspace produces: does it ``sound like an academic monograph rather than a
product white paper?''

This campaign also modeled review discipline that the BGS lane lacked. The
independent opinion accepted the findings as ``high-value triage input'' but
blocked execution of the rewrite until blocker-level repairs were
adjudicated, and every comparison tree carries a build gate recording the
commands actually run and their results. Review fed decisions; it did not
substitute for them.

\subsection{What the record adds up to}
\label{sec:lessons}

Ten conclusions survive contact with all five projects. Each is stated with
its mechanism and its design consequence, because a lesson without a
mechanism becomes a slogan, and slogans are how the BGS process failed.

\begin{enumerate}

\item \emph{Failures arrive in layers, and repair must follow the layer
order.} Format conversion, then register, then mechanical tells, then
defensive prose, then substance: the ZLB owner discovered the layers in that
order because each masked the next. The consequence for tooling is an
ordering constraint, not just a checklist: a style pass on an
internal-register document is wasted, and a substance review before the prose
is repaired will be answered by rewriting sentences the register pass would
have deleted.

\item \emph{Register leakage is contamination from the authoring
environment.} The BGS audit traced internal vocabulary on the page to
internal vocabulary in the author's working memory. This upgrades ``keep
process out of the prose'' from etiquette to architecture: reduce the
ceremony an agent must hold in context while drafting, and give drafting
agents contracts that name the audience and ban the internal lexicon
explicitly, as the ZLB style contract did.

\item \emph{Defensive prose is induced by integrity governance, and the cure
is relocation plus conversion.} Non-claims blocks and status labels were
mandated to prevent overclaiming and were then rejected by readers as
compliance theater. Deleting them wholesale would throw away real content, so
the ZLB repair converted them: prohibition becomes consequence, negative
definition becomes mechanism, and the accounting moves to an internal sidecar.
The editorial test that separated the two cases, reason or protection, is
implementable as a review prompt.

\item \emph{Templates have ceilings.} The uniform section shape that the ZLB
owner flagged was mandated by the project's own style contract, and every
checklist the BGS project introduced was promptly satisfied without the
quality it named. Floors get gamed and become ceilings. Pattern lists and
contracts therefore enter \hv{} only as subordinate diagnostics, and every
diagnostic reports rather than auto-fixes, so that no rule can silently
become the new template.

\item \emph{Local repair does not compose.} The BGS volume absorbed 815
paragraph repairs and still had no accepted reader path; the ZLB rescue
needed a register rewrite, not better sentences, before sentence work meant
anything. The unit of failure is the document's dependency graph. Sentence
diagnostics are still worth building, but nothing at sentence granularity can
certify a document, which is why the acceptance layer in
Section~\ref{sec:design} is reader-based.

\item \emph{Self-review cannot certify a human voice, and context-sharing
reviewers are contaminated.} Agent reviewers who shared the repository
accepted what a human rejected; the context-starved persona review predicted
the human. The design consequence is a review architecture: reader-proxy
reviews run with no repository context, forced choices, and an
instruction-injection boundary, and the human read remains the only
promotion gate.

\item \emph{Both failure directions exist, so length is not a signal.}
Documents in this record were rejected for being compressed summaries (the
BGS 21-page synthesis, the 23-page V11) and for being sprawling logbooks
(snapshot~06, the 115-page bridge edition); SMEwallet's accepted chapter was
longer than its rejected drafts. ``Make it shorter'' and ``make it longer''
are both category errors; the target is selection and recomposition with the
evidence layered elsewhere.

\item \emph{Two documents, not one.} Every attempt to make one artifact serve
as both audit record and human manuscript produced the hybrid readers
rejected. The internal record keeps ledgers, anchors, and identifiers; the
manuscript keeps none of them; an explicit identifier map ties them together.
The ZLB rescue kept the internal markdown as the record and the LaTeX as the
manuscript, with the code-name-to-math-symbol map written into the contract.

\item \emph{Measured diagnostics changed outcomes; unmeasured review did
not.} The ZLB repairs converged because someone counted 133 dashes and 20
identical openers, repaired, and recounted; ten unmeasured BGS review rounds
added nothing. Numbers also disciplined the arguments about whether a repair
had regressed. Every \hv{} diagnostic therefore emits counts with locations,
and acceptance tests for the tools are replays of these recorded numbers.

\item \emph{Guard the baseline or lose the right to edit aggressively.}
Byte-identical equation bodies, label and citation censuses, and
explain-every-removal reports are what let the ZLB editors rewrite freely
without anyone fearing silent damage; SMEwallet's protected-source maps did
the same job. The principle underneath is worth stating: readability
improvement is never established by shortening, so every removal must be
accounted for.

\end{enumerate}

One warning belongs beside the list rather than in it. The BGS project ended
up better at describing its failures than preventing them, and rewarding
honest self-assessment gave its author an exit ramp. \hv{} should expect the
same pull and should judge itself by one number only: how many rounds of
human feedback a document needs before acceptance. This survey, whose first
version failed that test immediately, is not exempt.

% Chapter: review of existing internal tools
\chapter{The tools already in the repository}
\label{sec:tools}

An author working in this workspace already has a prose policy, three skills, a
pattern list, two grounding tools, and a way to distribute them. Those pieces
can help, but they do not yet form a product. We need to know which ones help an
author in practice, which have evidence behind them, and where each one stops.
The rules are mature and largely consistent; executable measurement, behavioral
evidence, and one product-owned home are still missing.

Figure~\ref{fig:tool-gap} explains why consolidation alone is not the product.
The missing value lies in the handoffs between capable tools and in the reader
outcome at the end.

\hvFigureToolGap

\section{The prose policy}
\label{sec:tool-policy}

Every agent receives roughly 150 lines of the global policy under the heading
``Reader-Facing Scientific Prose.''\footnote{\path{claudecodex/policies/global-scientific-coding-agent-policy.md}. The companion section ``Plain Scientific Language And Non-Evasion'' supplies the verdict vocabulary discussed below.}
Its load-bearing device is a precedence order. When rules conflict, an agent
resolves them in this sequence: mathematical and scientific correctness
first, then source and evidentiary fidelity, then domain-appropriate meaning,
then reader comprehension and natural expression, and only then template
regularity. The order is what keeps a pattern list from doing damage. A rule
like ``remove em dashes'' sits at the lowest level, so it can never justify
deleting a numeric range from a theorem statement; a readability preference
can never justify dropping an assumption.

Three rules came directly from the failure record.

The first keeps internal workflow language out of ordinary exposition. Words
such as ``artifact,'' ``pipeline,'' ``gate,'' and ``handoff'' remain available
when the text is genuinely about software or project operations. They are not
used as metaphors for a theorem, an argument, or a paragraph. This rule turns
the administrative-register finding in Table~\ref{tab:bgs-taxonomy} into a
choice an editor can make sentence by sentence.

The second replaces defensive disclaimers with reasons. The policy prefers
``We assume log utility in this example because it keeps the first-order
condition simple'' to ``We are not claiming that households literally have
logarithmic utility.'' The first sentence states the modelling choice and its
purpose. The second introduces a mistaken interpretation merely to deny it.
This is the repair observed in Section~\ref{sec:bgs-audit}, expressed as a
general rule.

The third reserves acceptance for a reader. Author self-review, a model review,
a checklist, a style score, and a clean compilation can each provide evidence
about one part of a document. None can certify a human voice. The vocabulary
beside that rule is deliberately direct: a statement may be correct, wrong for
the stated target, unsupported, not checked, or heuristic. A polite synonym
must not blur which verdict applies.

Merits: the policy is complete in the sense that every failure class in
Chapter~\ref{sec:record} has a corresponding clause, and the precedence order
resolves the conflicts that sank naive rule-following elsewhere. Problems:
it is prose all the way down. Nothing in it executes, so compliance is
whatever the drafting agent happens to retain from 618 lines of context, and
the one thing the policy cannot do is notice that it is being ignored. The
ZLB rescue's third pass began with the admission that the policy ``had not
been consulted.'' A policy that must be remembered is a policy that will
sometimes not be.

\section{The narrative skill}
\label{sec:tool-linear}

When an agent uses \path{write-linear-scholarly-narrative}, it gets the
document-architecture skill,
distilled from the SMEwallet Draft~11 method.\footnote{\path{claudecodex/skills/write-linear-scholarly-narrative/SKILL.md}, with a 259-line companion, \path{references/teaching-contract.md}, defining the comprehension standard.}
Its sequence begins with the reader rather than the source files. The writer
states who will read the unit, what that person already knows, and what they
should be able to explain afterward without hunting through another document.
For this proposal, that means a former professor acting as an executive, not an
engineer who already knows the repository.

The writer then chooses the smallest case or question that can carry the
argument. Chapter order follows explanatory dependence: a concept appears
after the reader has encountered the problem that requires it, even if the
source material was originally stored in another order. Each bounded unit
opens one question and introduces the smallest idea needed to answer it.

Drafting follows a teaching sequence. Return to a known example, show the new
difficulty, and introduce notation only when it resolves that difficulty. An
equation is followed by its interpretation. The transition to the next unit
comes from a question the present unit cannot yet answer, not from a stock
phrase announcing that the document will move on.

Finally, a person reads the unit and tries to reconstruct the chain. They
should be able to explain why the unit began, what each turn contributed, what
the equation was for, and why the next question follows. Failure at that stage
is evidence about the exposition, not a request for the reader to study harder.

Its most valuable property is self-awareness about
Section~\ref{sec:lessons}'s fourth lesson. Nearly every rule carries an
anti-template clause; the drafting sequence is offered ``as a menu,'' with
the instruction not to fill slots mechanically and not to add a caveat or
transition merely because the template has a place for one. The
working-memory budget (about four live conceptual items) is labeled a
diagnostic threshold, not a law. Whoever wrote it had watched a floor become
a ceiling.

Evidence of effect is real but bounded: Draft~11's units passed their PI
gates one by one, which no earlier method achieved. Problems: the skill has
no examples inside it, so an agent meets its abstractions (``narrative
spine,'' ``reactivate'') without a single worked instance; it has no tooling,
so storyboards and contracts live or die by agent diligence; and its
acceptance step needs a human per unit, which priced it out of the BGS lane's
volume. Nothing tests whether an agent following it writes measurably better
than an agent that has never seen it, a gap Chapter~\ref{sec:proposal}'s
third work package exists to close.

\section{The sentence-level skill}
\label{sec:tool-natural}

The companion \path{write-natural-analytical-prose} works at sentence and
paragraph scale.\footnote{\path{claudecodex/skills/write-natural-analytical-prose/SKILL.md}, with before/after contrasts in \path{references/prose-patterns.md}.}
Before any rewrite, the agent
freezes a meaning ledger: the facts, attributed claims, deductions, live
uncertainties, author judgments, and exact numbers that must survive. After
the rewrite, it diffs the result against the ledger, and a smoother sentence
that changed a claim counts as a failure, not an improvement. Between those
two steps sit the craft rules: give each sentence one epistemic job, with a
calibrated verb table in which ``shows'' is not a synonym for ``suggests''
and ``reproduces'' is not a synonym for ``validates''; put concrete actors in
subject position (``the regime test rejects the authors' path,'' not ``the
evidence indicates''); state the finding before its qualifications; and spend
each qualification once, at the point of likely misunderstanding, rather than
re-attaching it to every positive sentence.

This is the skill that operationalizes the epistemic-blur and abstract-actor
rows of Table~\ref{tab:bgs-taxonomy}, and its meaning-ledger loop is the
sentence-scale version of the protected-baseline invariant. Its problems
mirror the narrative skill's: no automation, so the ledger diff is manual;
verification by ``read it aloud''; and hard prohibitions that depend entirely
on the drafting agent noticing its own violations. The skill is also, like
everything in this stack, tested only by phrase-lock (see
Section~\ref{sec:tool-claudecodex}): the test suite asserts that certain
sentences still appear in the skill file, not that the skill changes any
output.

\section{The humanizer pattern list, as installed}
\label{sec:tool-humanizer}

On 20~August the workspace vendored \texttt{blader/humanizer} v2.11.2 into
the policy tree, as
\path{humanizer-ai-writing-patterns.md}.\footnote{\path{claudecodex/policies/humanizer-ai-writing-patterns.md}. Section~\ref{sec:oss-humanizer} reviews the upstream project; this subsection reviews our installation of it.}
The list itself is 35 numbered AI-writing patterns derived from Wikipedia's
``Signs of AI writing'' catalogue, each with words to watch and a
before-and-after pair; Section~\ref{sec:oss-humanizer} walks through its
content. What matters here is the installation decision. The file carries a
locally written precedence header that demotes the whole list to ``a
DIAGNOSTIC pattern list, subordinate to
global-scientific-coding-agent-policy.md,'' spells out the order of
authority, and singles out the one rule that would otherwise do damage: the
flat dash ban of its pattern~14 is to be applied ``proportionately rather
than as a ban,'' and no pattern may be satisfied ``by deleting mathematics,
assumptions, or qualifications.''

That header exists because of a live incident. During the ZLB rescue the
list was consulted without a precedence frame, and the case study records
the realization that without one, ``a pattern list becomes a new template and
you trade one uniformity for another.'' The header is currently the only
policy-precedence mechanism in the workspace, and it is a comment in a
markdown file. It governs an agent exactly as far as the agent reads and
honors it. Turning this from a comment into configuration, per-format rule
downgrades in the diagnostic layer itself, is one of the cleaner design
requirements to come out of this review
(Chapter~\ref{sec:design}, under configurable policy precedence).

\section{The repair skill and the only working scripts}
\label{sec:tool-repair}

MathDevMCP ships its own scholarly-writing skill,
\path{repair-scholarly-readability}, and it is the only piece of the stack
with executable parts.\footnote{\path{MathDevMCP/skills/repair-scholarly-readability/SKILL.md}, with three reference rubrics and two scripts under the same directory.}
Where the narrative skill teaches drafting, this one teaches repair of an
existing LaTeX document, and its structure encodes two ideas the rest of the
stack lacks. The first is four separate records: reader comprehension,
mathematical integrity, source fidelity, and typography, with the explicit
rule that ``evidence in one ledger cannot certify another.'' A clean compile
is typography evidence only; a passed equation audit says nothing about
motivation. The second is bounded repair: one mechanism chapter or 6--12
rendered pages at a time, against a preserved baseline, ending in a test with a
reader who has no project history, prespecified questions, and a four-record
decision table.

Both scripts matter beyond their modest size.
\path{scripts/scholarly_surface_audit.py} emits JSON and markdown surface
diagnostics for a .tex/.pdf pair, and \path{scripts/render_review_pages.sh}
renders page images for visual inspection, institutionalizing the ZLB
rescue's discovery that text extraction lies about layout. These are the
embryo of the \texttt{hv capture} command in Chapter~\ref{sec:proposal}.

Evidence of effect: the skill's bounded method took the opening of the BGS
industry-DSGE dossier through a repair that passed five prespecified
reader-without-project-history questions, the strongest automated-plus-human result in the BGS
lane. The same result file records the limit: later chapters ``remain dry,''
and the project's audit concluded the skill ``generalizes locally, not across
a corpus-sized argument,'' which is lesson five of
Section~\ref{sec:lessons} observed in the field. The other problem is
organizational: this skill overlaps the claudecodex narrative skill in scope
while living in a different repository with different reference rubrics.
Scholarly readability currently has two homes and no owner.

\section{ResearchAssistant}
\label{sec:tool-ra}

When a citation needs checking, ResearchAssistant supplies the source-grounding
half of the discipline. It is a local-first tool that ingests papers (preferring arXiv LaTeX source over PDF,
because the LaTeX parses reliably and the PDF does not), extracts sections,
equations, labels, and bibliographies, tracks provenance and review status,
and exposes a read-only MCP surface plus a
CLI.\footnote{\path{research-assistant/}; the checked-in MCP configuration targets Claude Code and VS Code. Python 3.11 only.}
For this survey its relevant property is precedent: its \texttt{ra survey}
subcommands are the one place where a piece of workspace writing policy
became software. The literature-audit policy demands six ledgers (source
support, citation metadata, backward and forward snowballing, claim support,
omitted-paper risk); \texttt{ra survey central-papers} writes exactly those
ledgers and returns per-paper dispositions.

Its behavior under failure is the part worth copying. In the ZLB campaign
the tool's public-discovery bootstrap was unavailable, and instead of
degrading silently it closed the mission as
\texttt{terminal\_blocked\_bootstrap\_unavailable}, leaving the ledgers to be
maintained by hand. A tool that reports its own incapacity honestly is rarer
than it should be. Its limits for \hv{} are equally clear: it grounds
sources, not prose; its equation extraction from PDFs is explicitly marked
unreliable (arXiv LaTeX is the reliable path); and its discovery pipeline
depends on external metadata services that were down when most needed.

\section{The mathematics and modeling tools}
\label{sec:tool-math}

Two more tools border the writing problem without being writing tools.
MathDevMCP's mathematical layer provides SymPy-backed equality checking,
derivation audits, and an experimental whole-document rigor audit over
LaTeX. Its one recorded contribution to a manuscript is instructive about
scale: the ZLB audit produced per-equation issue ledgers and exactly one
applied exposition patch (an invertibility condition stated before a
displayed inverse), tagged
``candidate\_exposition\_patch\_not\_certificate.''\footnote{\path{BayesFilter/docs/surveys/zlb_discontinuous_hmc/mathdevmcp_audit_20260819.md}.}
That is the correct division of labor, mathematics tools nominating exposition
fixes without certifying prose, and \hv{} should preserve it rather than
absorb it. DynareMCP, the macro-model execution server, contributes structure
inventories that were used to falsify inflated rewrite claims (counting
inherited words and relocated blocks), a niche but genuinely useful
falsification tool for ``we rewrote it'' assertions.

\section{Distribution and the limits of the current tests}
\label{sec:tool-claudecodex}

claudecodex is the delivery mechanism. Its installer splices the global
policy into marked blocks in each agent's configuration files idempotently,
merges approval rules, and copies the skills into the Codex skill
tree.\footnote{\path{claudecodex/install_global_agent_policy.py}; templates and review-gate scripts under \path{claudecodex/docs/templates} and \path{claudecodex/scripts}.}
Around it sit the cross-agent review plumbing (bounded review packets, probe
ladders, a gate script that parses ``VERDICT: AGREE or REVISE'') and a
five-file proposal template system with seven evidence ledgers. The
machinery works and is in daily use.

Its test suite deserves a hard look because of what it teaches about the
stack's blind spot. There are real behavioral tests for the review gate (a
fake worker exercises the probe-and-fallback paths) and for the installer.
For the writing skills there are only phrase locks:
\path{test_natural_analytical_prose_skill.py} asserts that certain sentences
still occur in the skill file and that banned phrasings do not. Phrase locks
prevent silent regression of the doctrine's text. They say nothing about
whether the doctrine changes what an agent writes. As of today, no experiment
anywhere in the workspace has compared documents produced with and without
any of these skills. The doctrine is plausible, consistent, and untested.

Two further operational gaps matter. Installation is
Codex-first; Claude Code agents receive the policy but not a native skill
installation. And the writing doctrine now lives in three places, the
claudecodex policy tree, the claudecodex skills, and the MathDevMCP skill,
with cross-references in prose instead of a shared source of truth.

\section{The gap, stated precisely}
\label{sec:tool-gap}

The last column of Table~\ref{tab:tool-gap} gives the build list. The six
capabilities requested by the ZLB case exist nowhere as a coherent product: a
register linter, rhythm profiler, defensive-register finder, LaTeX-aware
baseline differ, a private provenance map that keeps evidence attached to one
manuscript, and a precedence mechanism stored in configuration rather than a
comment. The doctrine also has no behavioral evidence, and its three copies
need a shared source. Humanvoice joins that list into one testable review path.

\begin{table}[t]
\centering
\small
\begin{tabularx}{\textwidth}{@{}l L L@{}}
\toprule
Piece & What it reliably provides & What it cannot do \\
\midrule
Prose policy & complete doctrine; precedence order & execute, or notice being ignored \\
Narrative skill & Draft-11 method; anti-template clauses & show examples; run without a human per unit \\
Sentence skill & meaning-ledger safety loop; verb calibration & automate the ledger diff \\
Humanizer (installed) & 35-pattern diagnostic with FP guards & bind its own precedence outside prose \\
Repair skill & four-ledger repair; the only scripts & compose beyond 6--12-page units \\
ResearchAssistant & source grounding; ledgers as software & anything about prose \\
MathDevMCP/DynareMCP & equation audits; structure inventories & certify exposition \\
claudecodex & idempotent install; review plumbing & test writing behavior; reach Claude natively \\
\bottomrule
\end{tabularx}
\caption{The existing stack, by what each piece does and does not provide.}
\label{tab:tool-gap}
\end{table}

The inventory is therefore an asset and a warning. It gives the build a set of
tested interfaces, historical fixtures, and language for describing failure;
it does not give the project a validated intervention. The next part places
those local assets beside external evidence. That comparison matters because
a tool can be useful in this repository for reasons that do not travel to
another genre, and a widely used package can still fail the protected-source
contract. The field review begins with what research supports, then explains
how the literature was assembled, before returning to the package choices.

% Part 3: the research literature, explained
\section{What the research literature says}
\label{sec:literature}

This section explains the studies a designer of \hv{} needs to understand,
at the level of what the researchers actually did and found. Each subsection
ends with the consequences for our design, so that
Section~\ref{sec:design} can cite conclusions rather than re-argue them. A
note on inspection depth: the load-bearing papers below were read at method
level (design, data, and headline numbers verified from the full text or
extended abstract); the supporting papers were read at abstract level and are
cited for orientation only. One result, StoryScope, reached us secondhand
through an open-source project and is marked accordingly wherever it appears.

\subsection{The fingerprint: what makes text read as machine-written}
\label{sec:lit-fingerprint}

Start with the study that made the phenomenon measurable.
\citet{kobak2025delving} asked how much of the biomedical literature is now
LLM-assisted, without using any detector. They parsed more than fifteen
million English PubMed abstracts from 2010--2024 and tracked, for every word,
the fraction of abstracts containing it each year,
\begin{equation}
p_{t}(w) \;=\; \frac{a_{t}(w)+1}{b_{t}+1},
\label{eq:kobak-p}
\end{equation}
where $a_t(w)$ counts abstracts in year $t$ containing word $w$ and $b_t$ is
the year's abstract total. Borrowing the excess-mortality design from
epidemiology, they project a counterfactual 2024 frequency from the pre-LLM
trend,
\begin{equation}
q_{2024}(w) \;=\; p_{2022}(w) \;+\; 2\,\max\bigl(p_{2022}(w)-p_{2021}(w),\,0\bigr),
\label{eq:kobak-q}
\end{equation}
that is, the 2021$\to$2022 change extrapolated two years forward and clamped
so the counterfactual never falls (2023 is skipped as already contaminated).
A word is in excess if its gap $\delta = p - q$ or its ratio $r = p/q$
clears a threshold. The result that matters for us is \emph{which} words
were in excess. The Covid years produced excess content words (nouns about
the pandemic); 2024 produced excess \emph{style} words: ``delves'' at a
ratio near 28, ``underscores,'' ``showcasing,'' ``potential'' with a gap of
0.045. From the share of abstracts containing at least one excess style
word, they bound LLM assistance below by 13.5\% of 2024 abstracts, up to
40\% in some subcorpora. The bound is a floor, since authors who scrub the
telltale words become invisible, and that caveat cuts both ways for us:
scrubbing words defeats this measurement without improving the prose.

Why do models overuse exactly these words? \citet{juzek2025delve} isolated
the vocabulary by a three-way intersection: words whose PubMed frequency
jumped significantly between 2020 and 2024, that ChatGPT also overuses when
asked to write abstracts for 2020 papers (so the comparison is
like-for-like), yielding 21 focal words. They then tested explanations. The
words are not overrepresented in plausible training corpora, which weakens
the training-data story. The sharpest test compares Llama-2-7B Base with its
Chat variant, identical architecture, the Chat model adding only supervised
fine-tuning and RLHF: the Chat model is markedly less surprised by
focal-word-laden text. The overuse enters at the human-feedback stage. One
more of their findings matters for design: raters recruited to resemble the
RLHF workforce showed no overall preference \emph{for} the focal words, so
nothing in the reward signal pushes back against amplification once it
starts.

For practitioners, the catalogue of the fingerprint is Wikipedia's ``Signs
of AI writing,'' maintained by the editors who clean AI text out of Wikipedia
\citep{wikipedia2026signs}. It goes well beyond vocabulary: inflated
significance (``stands as a testament''), superficial \mbox{-ing} analysis
tails, rule-of-three padding, negative parallelism (``not just X, but Y''),
formulaic challenges-and-future sections, essay-like conclusions, title-case
headings, and chatbot residue. Two properties make it the right seed list
for diagnostics. It is maintained against live cleanup experience, so it
updates as the models change; and its preamble insists that no single sign
is diagnostic, since human writers produce all of them at some rate. That
warning is the germ of the false-positive discipline that the best
open-source tools in Section~\ref{sec:oss} implement and the worst ignore.

What about ``slop'' as a quality construct rather than an authorship one?
\citet{shaib2025slop} asked nineteen experts across NLP, writing,
journalism, linguistics, and philosophy to define it, coded the definitions,
and landed on three themes: information utility (density and relevance),
information quality (factuality, bias), and style quality (repetition and
templatedness, coherence, tone). Professional copy-editors then annotated
news and retrieval-augmented QA passages, first with a binary slop judgment,
then with word-level spans tagged by code. Two findings carry over to us.
Regressions showed \emph{relevance and density}, not style, best predict the
binary slop judgment, which is a formal version of our record's lesson that
substance failures hide under style failures. And even trained annotators
agreed only moderately (theme-level Krippendorff's $\alpha$ of 0.34--0.45),
which calibrates how much precision we should expect from any judge, human
or model, on this construct.

Finally, homogenization. \citet{padmakumar2024writing} hired 38 writers to
produce short argumentative essays in three within-subject arms: alone, with
base GPT-3 suggestions, and with feedback-tuned InstructGPT suggestions, with
keystroke logging attributing every character to writer or model. Essays
co-written with the instruction-tuned model were measurably more similar to
one another (higher pairwise Rouge-L and BERTScore within topics, fewer
distinct n-grams and key points), the base model showed no such effect, and
the decomposition traced the homogenization to the model-contributed text
specifically. The instruction tuning, not the assistance, is the
homogenizer. \citet{doshi2024generative} found the same shape in a
creative-writing field experiment, individual quality up, collective
diversity down, and \citet{guo2023curious} showed the spiral direction:
models trained recursively on synthetic text lose lexical, syntactic, and
semantic diversity generation over generation. One discourse-level result is
worth logging despite its secondhand status: the StoryScope study is
reported to detect machine narrative at 93\% F1 from discourse features
alone, with professional surface rewriting moving detection by less than two
points \citep{russell2026storyscope}. We have not inspected the paper; if it
holds, it independently confirms that the durable fingerprint is structural,
which is also what our failure record found when 815 sentence repairs left a
volume unaccepted.

\paragraph{Consequences for the design.} The fingerprint is distributional
and structural, not a word list. So \hv{} measures excess against a genre
reference corpus, Kobak-style, with
equations~\eqref{eq:kobak-p}--\eqref{eq:kobak-q} applied to pre-2022
economics text rather than PubMed; it treats word lists as seeds for
measurement, not rewrite targets; it measures homogenization \emph{across}
our documents, not only within one; and it expects style repair alone to
leave the deeper fingerprint intact, which is why the acceptance layer is
reader-based.

\subsection{Why the models write this way}
\label{sec:lit-training}

Four studies, read together, explain the register our agents default to, and
the explanation changes what interventions can work.

\citet{kirk2024understanding} trained the same base models through the
alignment stages, supervised fine-tuning, best-of-$n$ sampling, and
PPO-based RLHF, on summarization and instruction-following, then measured
performance (GPT-4-judged win rates, in and out of distribution) and output
diversity (distinct n-grams, embedding dispersion, and an entailment-based
measure, per input and across inputs). RLHF wins on performance, in and out
of distribution, and pays for it with a substantial drop in per-input
diversity. Notably, turning the KL penalty up, the knob that is supposed to
keep the tuned model close to the base model, lowered performance without
restoring diversity. The diversity loss is not a mis-set hyperparameter; it
is what the objective buys.

Why does preference optimization sharpen so hard? \citet{zhang2025verbalized}
supply the cleanest mechanism: typicality bias in the preference data
itself. Human raters, for reasons cognitive psychology has documented for
decades (mere exposure, processing fluency), systematically favor familiar,
schema-congruent text. Model the reward as true utility plus
$\alpha \log \pi_{\mathrm{ref}}(y\,|\,x)$, a bonus for being typical under
the reference model. On 6{,}874 preference pairs whose two responses are
\emph{equally correct}, the fitted $\alpha$ is about 0.57--0.65 and highly
significant: raters reward typicality even with correctness held fixed.
Optimizing that reward provably sharpens the policy toward
$\pi_{\mathrm{ref}}^{\,1+\alpha/\beta}$, collapsing onto the mode whenever
many answers tie on utility, which is precisely the situation in prose
style. Their intervention is a prompt change: ask for \emph{five responses
with their probabilities} rather than one response. Requesting a verbalized
distribution recovers 1.6--2.1$\times$ the diversity of direct prompting on
creative tasks and about two thirds of the base model's diversity across
post-training checkpoints, at no retraining cost.

Verbosity has its own paper trail. \citet{singhal2023length} decomposed
where RLHF's reward gains come from on three preference datasets, and the
answer is mostly length. Within fixed length bins, reward improves little
(on one dataset, within-bin gains are a few percent of the total); shifting
mass toward longer outputs accounts for 70--90\% of the improvement. Their
sharpest demonstration replaces the learned reward with a pure length score:
PPO against ``longer is better'' reproduces most of the measured win-rate
gain of PPO against the real reward model. Padding is not a stylistic
accident; it is what the training objective paid for.
\citet{sharma2024sycophancy} complete the picture on the agreeable side:
preference-trained assistants systematically flatter and hedge because
raters prefer agreement, the same mechanism that seeds throat-clearing and
deference into technical prose.

\paragraph{Consequences for the design.} These are training-time artifacts,
so prompt-level pleading (``write naturally'') fights the model's objective
and loses; what works is changing the sampling question (verbalized
sampling), editing after generation, and measuring. Compression deserves
first-class status among edit operations, since verbosity is a paid-for
bias. And any LLM used \emph{inside} \hv{} as an editor or judge carries
these same biases, which is why the judging protocol below is designed
around them rather than around trust.

\subsection{Detection: how it works and why it cannot be a gate}
\label{sec:lit-detection}

\hv{} needs to understand detectors for two reasons: their failure modes
rule out an entire class of tempting designs, and their internal statistics,
used gently, are legitimate population-level instruments.

The detection literature is one idea refined three times. Machine text sits
where the model expects text to sit; human text does not. GLTR
\citep{gehrmann2019gltr} visualized this directly, coloring each token by
its rank in the predicted distribution, and lifted human detection accuracy
from 54\% to 72\% just by showing the ranks. DetectGPT
\citep{mitchell2023detectgpt} made it a statistic: perturb the passage
(mask spans, let T5 refill them, about a hundred variants), score original
and variants under the suspected source model, and flag when the original
sits at a local likelihood peak, since sampled text lies at peaks and human
text does not. It works (0.95 AUROC on their news benchmark) but needs
roughly the right source model and a hundred model calls per document.
Fast-DetectGPT \citep{bao2024fastdetect} removed the rewriting: one forward
pass gives the conditional next-token distribution at every position, and
the statistic compares the log-probability of each actual token with the
mean and spread of alternatives sampled from that same conditional,
\begin{equation}
\text{curvature} \;=\; \frac{\log p(x) - \tilde{\mu}}{\tilde{\sigma}},
\label{eq:fdg}
\end{equation}
with machine tokens scoring near the top of their conditionals (curvature
around 3) and human word choices looking like ordinary draws (around 0). A
surrogate scorer makes it black-box, it is 340$\times$ faster than
DetectGPT, and it flags 80\% of ChatGPT passages at a 1\% false-positive
rate on benchmarks. Binoculars \citep{hans2024binoculars} normalizes
differently: the ratio of a passage's perplexity under one model to its
cross-perplexity between two sibling models, which cancels the ``weird
prompt'' confound (their capybara problem) and detects across generators
with one global threshold, over 90\% detection at 0.01\% FPR on their
benchmarks.

Now the failures, which are the design-relevant part.
\citet{liang2023detectors} ran seven widely used detectors on two human
corpora: essays by US eighth-graders and TOEFL practice essays by non-native
speakers. The US essays passed nearly clean; the TOEFL essays were flagged
as AI at an average false-positive rate of 61\%, with 98\% flagged by at
least one detector. The mechanism is exactly the detectors' core statistic:
limited lexical range produces low perplexity, and low perplexity reads as
machine. Two prompt experiments nailed it: asking ChatGPT to enrich the
TOEFL essays' word choice dropped their false-positive rate to 12\%, and
asking it to simplify the US essays raised theirs to 57\%. Polished, formal,
conventional prose, which is to say \emph{scholarly} prose, sits on the
wrong side of this statistic by construction. \citet{weberwulff2023testing}
tested fourteen academic and commercial tools and found none reliable, with
accuracy below 80\% and easy degradation; \citet{krishna2023paraphrasing}
showed that one round of automatic paraphrase defeats the
perturbation-family detectors wholesale. Detection is unreliable exactly
where we would use it and trivially evaded by anyone trying.

There is one statistically respectable use, and it is population-level.
\citet{liang2024monitoring} estimate the \emph{fraction} $\alpha$ of a
corpus that is substantially LLM-modified without classifying any document.
Represent each document by which adjectives it contains; estimate each
adjective's occurrence probability under a human reference corpus ($P$,
pre-ChatGPT reviews) and an AI reference corpus ($Q$, generated reviews);
then choose $\alpha$ to maximize the mixture likelihood
\begin{equation}
\hat{\alpha} \;=\; \arg\max_{\alpha} \sum_{i} \log\bigl((1-\alpha)\,P(x_i) + \alpha\,Q(x_i)\bigr).
\label{eq:mixture}
\end{equation}
On constructed mixtures with known $\alpha$ the estimation error is under
two percent, an order of magnitude better and about $10^{7}$ times cheaper
than aggregating per-document detectors. Applied in the wild: 10.6\% of
ICLR~2024 review text substantially modified (up from 1.6\%), 16.9\% at
EMNLP, no change at Nature-family journals, and higher $\alpha$ in
late, low-confidence reviews. The estimator is meaningless for accusing an
individual document, by construction, which is exactly the property we
want.

\paragraph{Consequences for the design.} No per-document authorship verdict,
ever: the false-positive profile lands on precisely the writers and the
register we care about, and evasion is cheap anyway. What survives is
(i) curvature or Binoculars-style scores as \emph{corpus-level} smoothness
trends, advisory only, computed on detexed prose; and (ii) the mixture
estimator of equation~\eqref{eq:mixture} as the statistically honest way to
answer ``how much of our output still reads machine-typical this quarter,''
with $P$ estimated from pre-2022 economics text and $Q$ from our own
agents' raw drafts. Both enter \texttt{hv drift} in
Section~\ref{sec:proposal}, gated against per-document use.

\subsection{The edit paradigm: the strongest positive evidence}
\label{sec:lit-editing}

Most directly useful for \hv{} is
\citet{chakrabarty2024salvage}, because it tests the exact pipeline we
propose: generate, diagnose, edit. They generated about a thousand
paragraphs with three frontier models from instructions reverse-engineered
out of published prose, had professional editors mark them up, and
consolidated the editors' labels into seven idiosyncrasy categories:
clich\'e, unnecessary or redundant exposition, purple prose, poor sentence
structure, lack of specificity and detail, awkward word choice, and tense
inconsistency. Eighteen professional writers then produced 8{,}035
categorized span-level edits, the LAMP corpus. Three of their findings set
our expectations. Writing quality did not differ across the three models,
so the idiosyncrasies are an ecosystem property, not a vendor one. An
LLM-based span detector reached 0.46 precision against expert spans, but
the experts themselves only agreed at 0.57, so ``flag the bad spans'' is
inherently noisy for everyone and should feed a human or model editor
rather than an automatic rewrite. And the payoff experiment: in blind
rankings by the writers, edited AI text (whether the spans came from
experts or the detector) beat raw AI text decisively, raw output ranked
last 60\% of the time, while human-edited text still ranked first. Editing
recovers most of the gap; it does not close it.

Two smaller results refine the objective. \citet{cheng2025humt} built a
corpus-based metric of human-like tone and a steering method, and found
that users do not uniformly prefer maximal human-likeness: the optimum
depends on genre and task. ``As human as possible'' is the wrong target;
``appropriate to the register'' is the right one, which for us is fixed by
the field's own literature. And the style-transfer literature
\citep{jin2022style} is the cautionary map: pre-LLM transfer methods
chronically traded content preservation against style strength, which is
the trade our meaning-ledger and baseline invariants exist to refuse.

\paragraph{Consequences for the design.} Build the editor as
span-diagnose-then-edit with categories, not as one-shot rewriting; expect
span precision near 0.5 and design the workflow so a human or a second
model adjudicates flags; measure edit effect by blind pairwise comparison;
and never let an edit change a claim, which our existing meaning-ledger
skill already enforces in doctrine and \texttt{hv diff} will enforce in
code.

\subsection{Judging writing with models, without fooling ourselves}
\label{sec:lit-judging}

\hv{} needs cheap judgments of ``which version reads better,'' so the
LLM-as-judge literature is load-bearing, mostly for its list of traps.

\citet{zheng2023judging} established the baseline: on their multi-turn
benchmark, GPT-4's agreement with pooled expert votes was 85\%, against 81\%
human-human agreement, so a strong model is a usable proxy for preference
between chat answers. The same paper measured the biases operationally.
Position bias: judge two near-identical answers twice with order swapped;
every judge showed order-dependent flips. Verbosity bias: pad one answer
with rephrased duplicates of its own list items, adding no information;
two of three judges preferred the padded answer more than 90\% of the time.
Self-enhancement: models rate their own outputs higher, though their design
could not settle causality. \citet{panickssery2024llm} settled it with a
clean causal design on summarization: models recognize their own outputs
above chance out of the box, and fine-tuning a model to be \emph{better} at
self-recognition increases its self-preference linearly, while control
fine-tunes on unrelated tasks do not. Self-preference rides on
self-recognition, so the generator must never judge its own rewrite.
\citet{dubois2024length} showed the verbosity bias is largely removable by
regression: length-controlled win rates change model rankings and nearly
triple agreement with human evaluations on their benchmark. Finally,
\citet{chakrabarty2024art} calibrate optimism for prose specifically:
expert writers applying a creativity rubric failed LLM-written stories on
roughly ten times more rubric items than professional stories, while LLM
judges rated the same stories generously. For writing quality, model judges
drift lenient exactly where experts are harsh.

\paragraph{Consequences for the design.} The judging protocol writes
itself, negatively: pairwise, never absolute scores; order randomized or
both orders averaged; length-controlled; a different model family from the
generator and from the editor; one rubric dimension per call rather than
holistic ``is this good''; and periodic validation of the whole proxy
against recorded human verdicts, which our evaluation corpus contains. Any
result that survives those controls is worth acting on; nothing else is.

\subsection{Readability science: measure against the genre, not a formula}
\label{sec:lit-readability}

Classical readability formulas are the wrong instrument for this project,
and it is worth knowing why rather than just asserting it. Flesch-family
scores are functions of syllable counts and sentence lengths. LLM prose
often scores \emph{well} on them while failing readers, because its
failures live in cohesion, rhythm, and register, which the formulas do not
see. The modern evidence: the CLEAR corpus work collected thousands of
passages with human ease-of-reading judgments and showed learned models
beat the classical formulas substantially \citep{crossley2022clear};
Coh-Metrix demonstrated two decades ago that texts separate on cohesion and
discourse dimensions the formulas ignore \citep{graesser2004cohmetrix}; and
Biber's register work established that genres occupy different regions of a
multidimensional feature space, so ``natural'' is genre-relative
\citep{biber1988variation}. Two findings about scientific prose complete
the picture. Scientific abstracts have been getting harder to read for a
century, with jargon and sentence complexity rising steadily
\citep{plavensigray2017readability}, so the median published abstract is a
degrading target, not a gold standard. And jargon density carries a
professional price: titles and abstracts heavier in specialized terminology
collect fewer citations \citep{martinez2021terminology}, a result an
economist will recognize as an incentive argument for plain language, not
just an aesthetic one.

\paragraph{Consequences for the design.} Every \hv{} metric is referenced
to a target-genre human corpus (for us: pre-2022 economics working papers
and journal articles), never to an absolute scale; the diagnostic suite
spans lexis, sentence rhythm, and discourse cohesion rather than syllable
arithmetic; and ``write like the median abstract'' is explicitly not the
objective, which the policy's plain-language stance already reflects.

\subsection{Economics and publishing norms}
\label{sec:lit-econ}

Norms here are settled enough to build on. \citet{korinek2023generative},
the profession's de facto guide, treats LLM writing assistance as ordinary
productivity technology for economists, to be used and disclosed. Journal
policy has converged on disclosure rather than prohibition (Nature
portfolio: models cannot be authors, use must be documented; Science moved
from ban to disclosure). The prevalence work says assistance is already
everywhere: the mixture estimator of Section~\ref{sec:lit-detection} puts
substantial LLM modification at 7--17\% of recent ML conference review text
\citep{liang2024monitoring, liang2024mapping}; AI-assisted reviews
measurably shift scores and acceptance at ICLR \citep{latona2024lottery};
adoption in research writing skews toward non-native English speakers and
junior researchers, whose styles are converging \citep{lin2025divergent};
and undisclosed chatbot artifacts (``as an AI language model\ldots'')
appear verbatim in published papers \citep{glynn2024academai}.

\paragraph{Consequences for the design.} \hv{} is a quality project under
disclosure norms, full stop. Detector evasion is not a goal, would not be an
achievement (Section~\ref{sec:lit-detection} showed evasion is cheap), and
is ruled out as a success metric because detectors cannot define quality.
The prevalence work also hands us our reference corpora: pre-2022 economics
text is a clean human baseline; anything after 2023 is contaminated and
unusable for calibration.

% Part 4: open-source tools, reviewed in detail
\section{Open-source tools, reviewed}
\label{sec:oss}

We evaluated fourteen open-source projects: the four named in the project
brief and ten found by searching for AI-tell skills, prose linters,
detectors, and LaTeX-aware writing tools. Every project was cloned and read
at the level of its actual pattern lists, prompts, rules, and code, not its
README; the runnable ones were run on our own documents. The clones are
pinned under \path{humanvoice/vendor/} with a commit
manifest.\footnote{\path{humanvoice/vendor/MANIFEST.md} records the twelve
cloned repositories with their commit hashes as of 21 August 2026; two
further tools (TeXtidote, textlint) were surveyed without cloning.}
This section reviews each tool in enough detail to defend its verdict:
what it contains, how it behaves, what is good, what is wrong, and what
adopting it would concretely mean. Table~\ref{tab:oss} at the end collects
the verdicts.

\subsection{A calibration experiment first}
\label{sec:oss-calibration}

One experiment frames every review that follows. sloptrim
(Section~\ref{sec:oss-sloptrim}) is the best deterministic AI-tell scorer we
found: 0--100, fully local, documented patterns. We ran it on three
documents whose human verdicts we know. Table~\ref{tab:calibration} shows
the result.

\begin{table}[t]
\centering
\small
\begin{tabular}{@{}lccl@{}}
\toprule
Document & Score & Band & Human verdict \\
\midrule
ZLB survey, internal markdown (pre-rescue) & 45 & mixed & rejected (register) \\
ZLB survey, final LaTeX (post-rescue) & 25 & light tells & \textbf{accepted} \\
This survey, first version & 10 & clean & rejected (does not teach) \\
\bottomrule
\end{tabular}
\caption{sloptrim v0.9.2 scores versus human verdicts. Lower scores mean
fewer mechanical tells. The scorer correctly separates the ZLB rescue's
before and after, and rates the rejected first version of this survey
cleaner than the accepted manuscript.}
\label{tab:calibration}
\end{table}

Read the table twice. First reading: the scorer works on its own layer. It
separated the ZLB before (45, driven by boldface density and inline header
lists, the markdown-era register) from the ZLB after (25), in the right
direction, deterministically, in about a second. That is real value:
regressions on this layer become visible without a human. Second reading:
the layer is not the one that decides acceptance. The first version of this
survey scored 10, the cleanest of the three, and was the most decisively
rejected, because its failure, density without explanation, lives above the
mechanical layer. This is our record's layered-masking lesson
(Section~\ref{sec:lessons}) reproduced by instrument. Every tool below
should be read against this calibration: the question is never ``does it
certify quality'' (nothing here does) but ``which layer does it measure, and
how honestly.''

\subsection{blader/humanizer: the canonical pattern list}
\label{sec:oss-humanizer}

\emph{One markdown skill, 457 lines; MIT; 36.8k stars; active. Verdict:
adopted, already installed; we are current with upstream.}

humanizer is Wikipedia's ``Signs of AI writing'' operationalized as an
editing skill: 35 numbered patterns in six groups (content, language and
grammar, style, chatbot artifacts, filler and hedging, plus late additions),
each with words to watch and a before-and-after pair. The content patterns
cover inflated importance, name-dropping, \mbox{-ing}-phrase pseudo-analysis,
sales language, vague sources, and formulaic outlook sections; the style
patterns cover dashes, bold density, list formatting, title case; the
chatbot patterns cover residue like ``I hope this helps''; the late
patterns, 34 and 35, name two register defects our record knows well,
answering objections nobody raised and rejecting alternatives nobody
proposed, which makes the list unusually aligned with our
defensive-qualification class.

Two design features separate it from its many imitators. It has a serious
false-positive section, ``What not to flag'': em dashes alone prove
nothing, formal vocabulary is not a tell, scope statements and genuine
disclaimers stay, secondhand text (quotes, titles, discussed examples) is
never rewritten. And it has a fact-preservation rule: the rewrite may not
add or drop a claim, number, date, quote, or citation, with any unsupported
addition counted as an error. Those two sections are why our installation
(Section~\ref{sec:tool-humanizer}) could be made safe with one precedence
header; a list without them cannot be made safe at all.

Its one dangerous rule is pattern 14: a flat ban on em and en dashes unless
the writer's sample uses them, with a mandated final search for both
characters. Applied to scholarly LaTeX it would strip mandatory typography,
en dashes in numeric ranges and name compounds (``Nash--Cournot''). Our
precedence header already overrides it to proportionate application; the
configuration layer of Section~\ref{sec:proposal} will enforce the override
mechanically. We diffed our vendored copy against upstream: upstream
\texttt{main} is at the same content version (2.11.2; the version bump since
the last tag was a README-only change), so no update action is needed.

\subsection{hardikpandya/stop-slop: sharper, narrower, wrong register}
\label{sec:oss-stopslop}

\emph{One 68-line skill plus three reference files; MIT; 16k stars; last
substantive commit March 2026. Verdict: mine for patterns; do not install.}

stop-slop is an essayist's de-slopping skill: eight core rules, a quick
checklist, and a five-dimension 1--10 self-scoring rubric (directness,
rhythm, trust, authenticity, density; below 35 of 50, revise). Its
distinctive material is genuinely distinctive. \emph{False agency}: ``the
data tells us,'' ``the decision emerges,'' an inanimate subject performing a
human verb where a person should be named. \emph{Negative listing}: ``It
wasn't X. It wasn't Y. It was Z,'' the multi-sentence teaser stack.
\emph{Narrator-from-a-distance} and the interrogative-opener crutch. None of
these appear in humanizer, all of them appear in our agents' drafts, and all
three transfer to our catalogue, with one carve-out: ``the data tell us'' and
``the model implies'' are ordinary, correct econometric usage, so the false-agency
rule needs a domain exemption list before it goes anywhere near a manuscript.

Installing the skill itself is ruled out by its register assumptions. ``Kill
all adverbs. No -ly words'' would delete ``approximately,''
``asymptotically,'' and ``conditionally'' from an econometrics paper. ``No
passive constructions'' conflicts with methods prose where the actor is
irrelevant by design. ```You' beats `People''' and ``put the reader in the
room'' import a blog voice that scholarly writing correctly refuses. ``Two
items beat three'' and the ban on ``every/always/never'' collide with
quantifier language. These are not flaws for its intended genre; they are a
genre boundary, and it is on the wrong side of ours.

\subsection{leonxlnx/taste-skill: not a writing tool}
\label{sec:oss-taste}

\emph{Thirteen skill files, 6{,}670 lines; MIT; 78.6k stars; active.
Verdict: skip.}

Despite surfacing in every search for writing skills, taste-skill is a
frontend design pack: landing pages, brand kits, image-generation prompting,
redesign checklists. The writing content it does contain is
marketing-copy hygiene for web pages (register locks for button text, bans
on fake-precise numbers and ``Quietly trusted by''), not prose craft. We
keep it in the record for one sociological observation. Its em-dash rule is
an absolute zero, and the maintainer's comment explains why: the agent
``historically ignored em-dash limits when phrased as `use sparingly'.''
Proportionate instructions decay into non-compliance; flat bans survive
because they are checkable. That is an argument for putting proportionality
into a machine-checked configuration (where ``at most $n$ per section,
except in math'' is enforceable) rather than into prose instructions, and it
is the last time we will be surprised that every popular skill bans dashes
outright.

\subsection{itallstartedwithaidea/writing-agent: rejected on both counts}
\label{sec:oss-writingagent}

\emph{Node.js multi-agent framework with a 40-check QA runner and a
200-phrase blacklist; MIT; 20 stars. Verdict: skip.}

writing-agent is the one evaluated project whose goal conflicts with ours.
Its framing is detector evasion: checks are named ``Watermark Stripping,''
``Model Attribution Evasion,'' ``False Positive Rate Exploitation,'' and
``Non-Native Bias Exploit,'' the last being the deliberate exploitation of
the bias documented by \citet{liang2023detectors}. Section
\ref{sec:lit-econ} settled our position: quality under disclosure norms,
never evasion, so this orientation alone disqualifies it as a dependency.

Its engineering would disqualify it anyway, and it is worth one example
because ``40-point QA'' sounds substantial until read. Check 19, ``Watermark
Stripping,'' begins with the comment: ``SynthID watermarks are token-level
statistical patterns. We can't detect them directly, but we check that
content was substantially rewritten,'' and returns a pass with an
aspirational message. Roughly a third of the forty checks are of this kind,
unconditional passes wearing serious names. The ``perplexity'' check is an
uncommon-word ratio against a 200-word stoplist. The blacklist flat-bans
``moreover,'' ``furthermore,'' and ``in conclusion,'' legitimate scholarly
connectives whose clustering, not existence, is the tell. The few sound
numeric ideas inside (sentence-length standard deviation targets,
paragraph-length variation, type-token ratio) exist in better-engineered
form in sloptrim and vale-ai-tells. Nothing here needs mining.

\subsection{Vale plus tbhb/vale-ai-tells: the CI measurement host}
\label{sec:oss-vale}

\emph{Vale: markup-aware prose linter in Go; MIT; 6k stars; very active.
vale-ai-tells: a Vale rule package, 111 prose rules, 15 commit-message
rules, 18 experimental structural metrics; MIT; 73 stars; active, with a
false-positive test corpus in CI. Verdict: adopt with a curated
configuration.}

Vale is the standard way to run prose rules in CI: rules are YAML files
grouped into styles, findings come back with file, line, severity, and
count, and per-format configuration can enable, disable, or downgrade any
rule for any file glob. That last property is the point for us. Our
policy-precedence principle, ``diagnostics subordinate to domain
correctness,'' becomes an ordinary configuration file: a rule that is wrong
for LaTeX is switched off for \texttt{*.tex} without touching the rule.

vale-ai-tells is the rule package that makes Vale worth adopting now, and
its engineering quality is visible in the rule bodies. Three examples show
the range. A vocabulary-class rule is a token list with a message; a
structural rule is a formula; a defensive-register rule is a set of shape
patterns. Respectively:

\begin{quote}\small
\texttt{DefensiveHedges.yml}: flags shapes like ``This is more of X than Y,
but\ldots,'' ``It's not the most X, but\ldots,'' ``This may seem X,
but\ldots,'' with the message ``State your point directly without
pre-defending it.''

\texttt{Average\-Sentence\-Length.yml}: a metric rule whose formula is
words over sentences and whose condition is ``greater than 25,'' with a
message noting AI text tends toward uniformly long sentences.

\texttt{SentenceStartEntropy.yml}: \texttt{extends: script}, running a
small Tengo program that computes the entropy of sentence-opener
distribution per section and warns on monotony.
\end{quote}

Where the package earns adoption is its experimental style, the
measurable-diagnostics collection our case study asked for: sentence-length variance, sentence-start entropy and
repetition, paragraph-length variance, passive density, contraction
avoidance, transition repetition, and document-level tricolon density, all
as counts with thresholds. About forty ``Figurative*'' rules implement
stop-slop's false-agency idea programmatically (inanimate subjects with
human verbs, one verb family per rule so each span produces one alert). A
separate commit-message style (``This commit adds\ldots,''
``\ldots ensuring consistency'') transfers directly to our multi-agent
commit hygiene.

Four stock rules are destructive for our use and must be disabled or
downgraded for LaTeX, and it is worth recording exactly why.
\texttt{EmDashUsage} fires at error level on every em \emph{and} en dash,
so every numeric range fails. \texttt{DoubleHyphen} flags literal
\texttt{--}, which in LaTeX source is the correct en-dash syntax; it would
flag \texttt{1990--2020} in every bibliography. \texttt{FormalRegister}
bans, at error level, a token list including ``utilize,'' ``facilitate,''
and ``implement,'' and in the same list ``minimize'' relatives; an estimator
that minimizes an objective cannot be discussed under that rule.
\texttt{FormalTransitions} fights the existence of
``moreover/furthermore'' when the tell is clustering. All four are two
lines each in \texttt{.vale.ini}, which is precisely the precedence
mechanism as configuration.

The one real gap is LaTeX itself. Vale has no native \texttt{.tex} parser
(verified in its format registry, which lists markdown, HTML, reST,
AsciiDoc, and friends, but not TeX). The workable options are mapping
\texttt{tex} to the markdown parser with block-ignore regexes for math
environments, running Vale over detexed text at the cost of line fidelity,
or hosting rules in a genuinely LaTeX-aware tool
(Section~\ref{sec:oss-latex}). Piloting that choice is an explicit early
task in WP2.

\subsection{seyedehsanhadi/sloptrim: the deterministic scorer}
\label{sec:oss-sloptrim}

\emph{Single-file Python detector (2{,}223 lines, standard library only)
plus Claude Code hooks; Apache-2.0; 178 stars; active, 119 tests in CI on
three platforms. Verdict: adopt.}

sloptrim is the tool from Table~\ref{tab:calibration}. Its detector reads a
file (twenty formats, including \texttt{.tex}, \texttt{.docx}, and
\texttt{.ipynb}; prose is extracted, code is skipped), scores the first
256\,KB against its pattern catalogue, and emits JSON: per-pattern counts
with sample spans, structural metrics (sentence-length coefficient of
variation, paragraph monotony, opener repetition, hedge stacking, synonym
cycling, transition clusters, cadence zigzag), and a 0--100 score in five
bands with a word-count-aware confidence label; a flag adds bootstrap
confidence intervals. The hook layer wires the detector into Claude Code so
every file saved through the editing tools is scored on save and flagged
spans are surfaced to the agent; a documented blind spot is that files
written by shell redirection bypass the hooks.

Two properties earn the adopt verdict beyond the score itself. The pattern
accounting is honest: 71 documented patterns, 62 with detectors, and only 50
allowed to move the score, because the project measured the other 12 against
human formal writing and found they mark register rather than authorship,
demoting them to advice. That is our precedence principle, discovered
independently by measurement, inside the tool. And the project publishes its
own limits: a benchmark with ROC-AUC by arm (0.76--0.95), an explicit
statement that it is not an authorship classifier, and an ethics file
warning against consequential use.

Observed behavior on LaTeX, from our runs: equation environments and inline
math pass through unmolested (the \texttt{--clean} mode touches characters
only, never wording), which is the property we most needed. Two artifacts
to correct for: pattern words inside \verb|\text{...}| are scanned as
prose, and on raw \texttt{.tex} the opener-repetition metric can count
command sequences as sentence openers (our run flagged ``\verb|\item|
\verb|\textbf|'' repetition in a list). Scoring detexed text alongside raw
source removes both. The tool slots into \hv{} as the on-save scorer and
the first layer of \texttt{hv rhythm}, with the calibration of
Table~\ref{tab:calibration} rerun whenever its pattern set updates.

\subsection{AIScientists-Dev/academic-humanizer: the missing academic layer}
\label{sec:oss-academic}

\emph{One 261-line skill plus examples; MIT; 1.1k stars; last push July
2026. Verdict: adapt, by merging its content into our policy layer.}

This is the one evaluated project written for our genre, and it opens with
the sentence our whole stack is organized around: ``academic writing
already has a correct human voice: neutral, precise, third-person plural
(`we'), every claim tied to its evidence.'' No personality injection, no
casualization, explicitly ``not for evading AI-use disclosure.'' Its value
is three layers our installed humanizer does not have.

Layer 2 is an academic-specific tell catalogue, and its examples are chosen
well enough to quote. Over-claiming verbs: ``We prove that our method
significantly outperforms all prior approaches'' becomes ``Our method
improves held-out accuracy by 4--7 points over the strongest prior approach
(Table 3); the gain is significant at $p<0.01$ by a paired test.''
Significance hype (``paves the way,'' ``sheds light on''), empty
intensifiers (``extensive experiments on a wide range of datasets'' becomes
``we evaluate on three datasets, named''), novelty padding (``to the best
of our knowledge''), formulaic openers (``In recent years, X has attracted
increasing attention''), connective \emph{clustering} (the correct
formulation, against consecutive-sentence connectives rather than the words
themselves), citation dumping, contribution-list clich\'es, and, notably,
clause-stacked sentences past thirty words, with the instruction to split.
Every one of these appears in our agents' drafts; none is in Wikipedia's
catalogue.

Layer 3 is a preserve-list that functions as a built-in precedence header:
keep evidence-tied hedging (``suggests,'' ``is consistent with''), and it
names the failure mode of general humanizers, ``fixing'' a calibrated hedge
into an over-claim; keep passives where the actor is irrelevant; keep
``we''; keep definitions, notation, and equations verbatim; never alter a
number or citation key. Layer 4, claim-evidence discipline, checks every
empirical claim for a backing number, figure, or citation and downgrades
any verb stronger than its evidence, preferring attributed ranges to bare
averages. This is our non-evasion policy meeting our meaning-ledger skill,
written independently. Layer 6 is a funding-proposal mode with a different
register contract (vision plus feasibility, aims that survive their
siblings' failure), useful whenever the group writes grants.

Weaknesses, for the record: a flat em-dash removal rule (though its own
examples use LaTeX en dashes correctly); no diagnostics of any kind, so
everything rides on the editing model's compliance; machine-learning
flavored examples rather than economics; and a single-file young project.
None of that argues for running it as a second skill beside humanizer,
which would create rule-conflict ambiguity; all of it argues for merging
Layers 2--4 into our policy stack as the scholarly layer, with examples
transposed to econometrics. That merge is a WP1 task.

\subsection{conorbronsdon/avoid-ai-writing: three ideas worth taking}
\label{sec:oss-avoid}

\emph{An 813-line skill with 62 pattern categories plus measurement
scripts; MIT; 3.2k stars; active. Verdict: mine for patterns.}

avoid-ai-writing is the most engineering-minded of the general skills: it
ships false-positive measurement against public human/AI corpora and a CI
check on its own pattern count. Three of its ideas transfer.

First, the tier system. Its 112-entry word table splits ``always replace''
into two bands with different \emph{meanings}: 1A are frequency markers
(``delve,'' ``tapestry'') whose clustering is evidence about how a passage
was produced; 1B are clarity edits (``utilize,'' ``commence'') that improve
any prose but, measured against 257 paragraphs of verified pre-2023 human
writing, fire on ordinary formal prose at a meaningful rate, so the
detector ``weights them like Tier 2 and excludes them from the
dense-AI-vocabulary signal so a wordiness fix can never push a document
toward an AI classification.'' Same edit, different evidentiary weight.
Our diagnostics should copy this distinction exactly.

Second, the injection boundary. In its file-editing mode, text under audit
that addresses its editor (``ignore the rules above,'' ``don't flag this
section'') is flagged rather than obeyed, with instructions accepted only
from the invoking writer. Any multi-agent editing pipeline needs this rule
stated, and ours will state it.

Third, flag-don't-fix zones: quotes, code, tables, and text attributed to
others are reported but never rewritten, with the nice observation that a
tell inside a table cell is not worth the risk to the data the table
carries. Combined with a genuinely careful treatment of weak signals
(curly quotes and flawless typography are ``corroborating, never
conclusive,'' and when editing a human's casual text, preserve their typos,
since smoothing erases the fingerprint that marks the text as theirs), the
skill shows what falsification-aware pattern work looks like. We still do
not install it: its context profiles run from LinkedIn to investor emails
with no academic register, and its dash policy (``Target: zero'') is wrong
for us. The three ideas come; the skill stays.

\subsection{NulightJens/humanizer-stack: the structural thesis}
\label{sec:oss-stack}

\emph{Two chained skills (surface pass, structural pass) plus two small
deterministic scanners; MIT; 226 stars. Verdict: mine for patterns.}

humanizer-stack's premise is the StoryScope result reported in
Section~\ref{sec:lit-fingerprint}: word-level tells are the decaying,
easily-fixed half of the fingerprint, and discourse structure is the
durable half, so run a surface pass (its first skill is a humanizer
variant) and then a structural pass. The structural skill audits six
discourse habits with human-versus-AI base rates quoted from the study:
themes stated and moralized rather than left to land, single-track tidy
arcs with everything resolved, no digressions, unbroken linearity, and,
its deepest point, \emph{shape convergence}: all five models it examined
occupy one tight region of structural space while human pieces are
dispersed. Its operating instruction follows: pick one or two structural
interventions per piece and vary them across pieces, because a fleet of
documents fixed the same way converges on a new detectable cluster.

For our genre the six audits mostly do not transfer; a methods section is
supposed to be tidy and linear, and IMRaD structure is a feature. Two
things transfer anyway. The anti-convergence warning is directly relevant
to \hv{}, which will edit every manuscript the group produces with one
toolkit; measuring cross-document structural similarity is the guard, and
it enters the design as part of \texttt{hv drift}. And the two scanners
(\texttt{copy\_scan.py}, \texttt{structural\_scan.py}) are small clean
reference implementations of deterministic tell-counting, worth reading if
not importing. The skills themselves stay out: their genre calibration is
narrative and marketing, and their base rates do not describe scholarly
prose.

\subsection{Corpus instruments: slop-forensics and fast-detect-gpt}
\label{sec:oss-corpus}

\emph{sam-paech/slop-forensics: Python toolkit for deriving
over-representation lists from model output; MIT; 369 stars. Verdict:
adapt, corpus-level. baoguangsheng/fast-detect-gpt: research code for the
ICLR 2024 detector; MIT; 424 stars. Verdict: adapt, aggregate-only, under
guardrails.}

slop-forensics inverts the stale-ban-list problem. Instead of consuming
someone's 2024 word list, it takes a corpus of model outputs, compares
word, bigram, and trigram frequencies against human baselines (via
\texttt{wordfreq} and NLTK), computes vocabulary-complexity and slop-index
metrics, and emits the over-representation list for \emph{your} models on
\emph{your} tasks; it can even cluster models phylogenetically by slop
profile. Pointed at our own agents' drafts with a pre-2022 economics
baseline instead of general English (its stock baseline would flag
``heterogeneous'' and ``endogeneity'' as anomalies), it becomes a
self-updating tell list, which matters because the tells decay: ``delve''
is already fading as vendors patch it, and yesterday's list quietly loses
power. This is the Kobak method of
equations~\eqref{eq:kobak-p}--\eqref{eq:kobak-q} as maintained software,
minus the counterfactual projection, which we add ourselves in \texttt{hv
drift}.

fast-detect-gpt's repository contains clean reference implementations of
the curvature statistic of equation~\eqref{eq:fdg}, including a local
inference script that runs with small surrogate models on CPU. Under
Section~\ref{sec:lit-detection}'s findings, its only legitimate role here
is as an aggregate drift instrument: a quarterly curvature distribution
over detexed manuscript prose, watched for trend, never reported
per-document, never attached to a name, and interpreted with the standing
caveat that polished scholarly register inflates false positives by
construction. If that discipline ever feels burdensome, the correct
response is to drop the instrument, not the discipline. (Binoculars, the
other zero-shot detector we cloned for comparison, is stale since 2024 and
needs a heavier model pair; fast-detect-gpt supersedes it for local use.)

\subsection{The LaTeX question: TeXtidote and textlint}
\label{sec:oss-latex}

\emph{TeXtidote: LaTeX-aware grammar and style checker wrapping
LanguageTool; GPL-3; 1.1k stars; active. textlint: pluggable prose linter
with a LaTeX AST plugin; MIT; 3.2k stars core, plugin small; active.
Verdict for both: pilot.}

Section~\ref{sec:oss-vale} left one problem open: none of the AI-tell rule
hosts parses LaTeX natively. Two tools do. TeXtidote reads \texttt{.tex}
directly, skips math environments, maps every finding back to source
lines, and wraps LanguageTool, whose custom-rule mechanism (XML regex
rules) could host simple tells with real LaTeX awareness; GPL-3 is fine
for internal use. textlint is the opposite trade: a first-class custom-rule
ecosystem in JavaScript with a true LaTeX parser plugin
(\texttt{textlint-plugin-latex2e}), but no existing AI-tell rule pack, so
everything we want would be written by us. The pragmatic WP2 sequence is:
start with Vale over detexed text plus sloptrim on raw source (both work
today), and pilot TeXtidote as the line-accurate LaTeX layer; move rules
into textlint only if the pilot shows the detex mapping loses too much.

\subsection{proselint, and the rest}
\label{sec:oss-rest}

proselint (BSD-3, 4.6k stars) is the classic usage linter, 87 checks
distilled from Garner and the style canon: archaisms, clich\'es,
redundancy, weasel words. It predates LLMs, contains no AI-tell content,
and, decisively, has no custom-rule mechanism, so extending it means
forking it; since a Vale-compatible re-implementation exists, anything we
want from it arrives through the Vale configuration for free. Skip, and
inherit via Vale if ever wanted. The remaining survey finds were language
ports of humanizer (a 15.8k-star Chinese localization among them), smaller
Vale slop styles superseded by vale-ai-tells, and academic humanizer
variants less developed than the one reviewed above; none changes a
verdict.

\subsection{Verdicts, and what adoption means}
\label{sec:oss-verdicts}

\begin{table}[t]
\centering
\small
\begin{tabularx}{\textwidth}{@{}l l L@{}}
\toprule
Project & Verdict & What the verdict means in practice \\
\midrule
blader/humanizer & adopted & keep vendored copy; precedence header stays; enforce dash override in config \\
tbhb/vale-ai-tells + Vale & adopt & CI linter; disable four rules for \texttt{.tex}; enable structural metrics; take commit rules \\
seyedehsanhadi/sloptrim & adopt & on-save scorer and rhythm layer; recalibrate per Table~\ref{tab:calibration} on updates \\
AIScientists-Dev/academic-humanizer & adapt & merge Layers 2--4 into our policy as the scholarly layer, examples transposed to economics \\
sam-paech/slop-forensics & adapt & derive our own tell lists against a pre-2022 economics baseline \\
baoguangsheng/fast-detect-gpt & adapt & aggregate drift metric only; never per-document \\
conorbronsdon/avoid-ai-writing & mine & take 1A/1B weighting, injection boundary, flag-don't-fix zones \\
hardikpandya/stop-slop & mine & take false agency (with econ carve-out), negative listing, Wh-openers \\
NulightJens/humanizer-stack & mine & take the anti-convergence diagnostic; skills stay out \\
sylvainhalle/textidote & pilot & candidate line-accurate LaTeX layer \\
textlint + latex2e & pilot & fallback host if detex mapping proves lossy \\
amperser/proselint & skip & no custom rules; content reachable via Vale \\
leonxlnx/taste-skill & skip & frontend pack; one sociological lesson retained \\
itallstartedwithaidea/writing-agent & skip & evasion-framed; weak checks; nothing to mine \\
\bottomrule
\end{tabularx}
\caption{Verdicts over the evaluated open-source field, with the concrete
action each verdict implies.}
\label{tab:oss}
\end{table}

Put together: humanizer stays as the installed general diagnostic;
academic-humanizer's content becomes our scholarly layer; sloptrim and a
curated Vale configuration become the measurement floor; slop-forensics
and the mixture estimator keep the measurements current; three skills
contribute patterns without being installed; two LaTeX-aware hosts get a
pilot; and three projects are set aside for cause. Nothing in the field
does what Sections~\ref{sec:record} and~\ref{sec:tools} showed we need
most, reader-based acceptance testing and behavioral evaluation of writing
doctrine, which is why the proposal builds those two things and mostly
configures the rest.

\chapter{The first build}
\label{ch:proposal}
\label{sec:design}

The first version should solve one complete problem before it asks an executive
to read the result. Give it an authoring brief, the evidence the document may
use, the source material, and the objects an author cannot afford to change
silently. Humanvoice should return an argument blueprint, bounded draft units,
a reproducible rendering, a pre-human preflight record, cause-specific repairs,
and a reader packet tied to the exact released version. If it cannot establish a
correspondence or defend the argument, it should stop before the document
reaches the reader.

That vertical slice is smaller than the platform described in earlier drafts.
It is also more useful. A team can watch one consequential document travel from
brief to release, measure the work at each step, and return to the incumbent
workflow when the system abstains. The human is still the final authority, but
the human is no longer the default place where ordinary prose defects are
discovered.

The dependency graph in Figure~\ref{fig:mvp-dependencies} shows the critical
path. The protected LaTeX core comes before corpus scale, optional model
features, and any claim about comparative effectiveness.

\hvFigureMVPDependencies

\section{Ten commitments that shape the code}

The evidence narrows a credible implementation to ten commitments. Engineers
still choose the components, data structures, and interface within those
constraints.

\begin{table}[H]
\centering
\small
\begin{tabularx}{\textwidth}{@{}p{0.20\textwidth}L p{0.23\textwidth}@{}}
\toprule
Commitment & Product behavior & Observation that can overturn the choice \\
\midrule
One immutable version & Findings, objects, the PDF, and the reader session all
refer to the same source and render hashes. & A replay produces different
locations, objects, or pages from the recorded environment. \\
Brief before prose & A critical blank in reader, decision, evidence boundary,
or protected objects blocks long-form drafting. & A fluent draft is produced
from an incomplete or contradictory brief. \\
Blueprint before units & Claims, dependencies, section jobs, terminology, and
figures are checked before a writer inherits them. & A unit uses an untaught
concept, repeats a job, or makes a claim with no evidence anchor. \\
Bounded writing & The writer receives one approved unit at a time and cannot
widen the claim or invent evidence. & A whole-document prompt hides a failure
until the reader sees it. \\
Preflight before people & Integrity, argument, writing, and reconstruction gates
must pass before a reader packet exists. & A seeded defect reaches the packet,
or an unknown result is treated as pass. \\
Questions before rewrites & A diagnostic reports a span, evidence, and an
editorial question; inspection never changes source. & Authors cannot act on the
finding without asking a model to reconstruct its context. \\
Protected change, not frozen text & Equations, numbers, labels, citations,
tables, quotations, and declared claims can change only with a visible
disposition. & A seeded material change passes without notice. \\
Different human decisions & The author owns meaning and revision; the named
reader owns whether the document supports the task. & The interface pressures
either person to accept a system score or hides disagreement. \\
Local by default & Source, findings, diffs, and reader data remain local unless
an owner approves named fields for a named service. & A dependency requires an
undeclared network path. \\
Claims rise with evidence & Replays establish regression, held-out fixtures
establish operational fit, and later comparison estimates efficacy. & A
historical case is used to claim general effectiveness. \\
\bottomrule
\end{tabularx}
\caption{The main design commitments, stated as behavior and possible
falsification rather than internal requirement codes.}
\label{tab:requirements}
\end{table}

The machine-readable requirements retain stable identifiers in the appendix
and repository. The main argument uses names because a reader needs to
understand the tradeoff before tracing its implementation ID.

\section{The vertical slice}
\label{sec:proposal}

The slice begins with an authoring brief. The author names the document genre,
the intended reader, the decision, the concepts the reader can reasonably be
expected to know, the named exemplars that set the register, the evidence
boundary, and the objects that must survive revision. Humanvoice compiles that
brief into a claim map, concept-dependency graph, section cards, terminology
plan, and visual plan. A missing critical field blocks the long draft and leaves
an explicit question for the owner.

Only then does the system capture the source tree and canonical PDF. The primary
parser extracts source locations and protected objects; a second adapter or the
compiled representation checks the result. Unsupported constructs remain
explicit abstentions and block any gate that relies on them.

The first diagnostic set is deliberately small. It locates private repository
language, undefined project labels, broken references, conspicuous repeated
openings, paragraph packing, concept bursts, first-use ordering, and changes to
protected objects. These are observable conditions with stable locations. The
HPO case supplies calibration fixtures for the packing, pacing, and ordering
families; it does not turn their metrics into a pass/fail score. A more
interpretive question, such as whether a qualification replaces a reason, may
enter the feasibility interface only after human calibration shows that the
queue repays the attention it consumes.

The writer produces one unit from the approved plan, and the author can revise
that unit in the existing editor. Humanvoice records accept, reject, revise, or
defer for each finding, rebuilds the unit, and compares the new source with its
parent. Preflight critics attempt reconstruction before release. Only then does
the reader receive the clean PDF and a short reconstruction task in a local
browser. The event record joins the response to the version without showing
private rule names, confidence scores, or project history.

\begin{table}[H]
\centering
\small
\begin{tabularx}{\textwidth}{@{}p{0.18\textwidth}p{0.30\textwidth}L p{0.22\textwidth}@{}}
\toprule
Stage & Input & Result & Failure behavior \\
\midrule
Brief & Reader, decision, evidence boundary, exemplars, and protected objects &
Versioned authoring contract. & Stop on a missing critical answer; do not draft
the long document. \\
Plan & Brief plus source inventory & Claim map, dependency graph, section cards,
and visual plan. & Block a claim without evidence or a concept used too early. \\
Draft & One approved unit and its evidence & Rebuildable prose and source map. &
Keep the unit internal until its checks pass. \\
Preflight & Draft, protected manifest, build, critics, and mutations & Gate record
with located failures and abstentions. & Fail closed; no reader packet. \\
Repair & Gate failure and named repair strategy & New unit, comparison, and
convergence record. & Stop on cycle limit or oscillation and escalate to owner. \\
Release and read & Passed gates, clean PDF, decision context, and timed task &
One reader packet and a human decision. & Preserve noncompletion and decline as
outcomes rather than deleting the session. \\
\bottomrule
\end{tabularx}
\caption{The complete first-release path and its explicit fallbacks.}
\label{tab:mvp-vertical-slice}
\end{table}

\section{The boundary of the first release}

Twelve weeks buys a LaTeX-first implementation. It includes one primary parser
adapter and one fallback, the canonical build, the protected-object manifest,
the bounded diagnostics above, author dispositions, the protected comparison,
a self-hosted reader task, and an exportable event record. The parser bake-off
includes pylatexenc, unified-latex, TexSoup, and the richer document-model
candidates described in Chapter~\ref{ch:software-dossiers}; a package wins only
by passing the locked fixtures with acceptable operator burden.

Markdown waits. So do a general editor extension, unbounded autonomous rewriting,
continuous corpus-drift monitoring, broad rhetorical-move inference, and a
market-facing multi-tenant service. Local model suggestions remain an optional
experiment behind the same brief, blueprint, preflight, and author-decision
record; the deterministic path must still be useful. Single-document
authorship classification is not a deferred feature. It is outside the product
boundary.

The narrow scope produces a useful fallback. If the diagnostic layer fails to
earn attention, the source snapshot and protected comparison can still become a
small review utility. If the parser fails too often for full correspondence,
the reader task can still test whether structured low-context review is useful.
A platform with seven interdependent commands would make those negative results
harder to learn from and more expensive to salvage.

\section{A small command surface}

The interface exposes six verbs. Their names are less important than their
contracts, and each one has a fail-closed boundary.

\begin{description}[style=nextline,leftmargin=28mm,labelindent=0mm]
\item[\texttt{hv init}] Validates the authoring brief, evidence boundary,
privacy class, and protected-object declarations. It refuses a missing critical
field.
\item[\texttt{hv plan}] Builds and checks the claim map, concept-dependency
graph, section cards, terminology order, and visual plan. It emits questions,
not prose, when the evidence is insufficient.
\item[\texttt{hv draft}] Invokes a replaceable writer for one approved unit; the
unit receives only the context and evidence named in the blueprint.
\item[\texttt{hv preflight}] Runs integrity, argument, writing, evidence,
reconstruction, and mutation checks. Unknown results remain blockers.
\item[\texttt{hv repair}] Applies a named repair strategy, rebuilds the unit,
reruns preflight, and stops on the cycle limit or oscillation. It writes a
child revision and never edits the canonical source in place.
\item[\texttt{hv release}] Emits the clean reader packet and launches the local
timed task only after the release gate passes. It records completion, response,
confidence, and decision against one immutable PDF.
\end{description}

The older capture, inspect, and compare services remain internal adapters for
these commands. Citation identity and build diagnostics are part of preflight,
not optional reports. Population monitoring and research-corpus analysis live in
the evaluation code, not in an author's release command. This keeps a research
instrument from acquiring editorial authority through a convenient interface.

\section{Regression, calibration, and held-out validation}

The earlier proposal asked historical documents to do three incompatible jobs.
The revised plan separates them.

\paragraph{Regression fixtures.} Frozen ZLB snapshots test whether the software
can reproduce recorded facts: the 133-dash snapshot, the distribution of
section openings, the 109 unchanged equation bodies, and the deliberately added
references. These tests catch implementation drift. They say nothing about
performance on a new document.

\paragraph{Calibration documents.} A declared training partition supplies
examples for rule wording, thresholds, source-location repair, and reviewer
instructions. Every tuning decision is logged before the held-out partition is
opened. A document that changes a rule can no longer estimate that rule's
operational value.

\paragraph{Pre-human mutation tests.} Before a human sees a packet, the harness
injects known failures into otherwise valid plans and drafts: a missing evidence
anchor, a use-before-teach concept, a dense paragraph that exceeds its declared
concept budget, a citation substitution, a changed sign, a broken reference, an
internal project label, and a repair that oscillates. Critical deterministic
mutations must be caught or explicitly abstained from. Rhetorical mutations are
reported with confidence intervals and a frozen threshold chosen on calibration
data; no threshold is tuned on the live case.

\paragraph{Held-out operational validation.} New non-ZLB fixtures include
custom macros, malformed environments, moved equations, changed signs and
 units, valid citations with unsupported claims, legitimate repeated prose,
 private labels a reader does not know, and the HPO first-use positives and
 adjudicated false alarms. A second operator replays the package adapters from
 a clean environment. Held-out authors judge whether each prose finding was
 worth attention, and those authors are not the people who wrote the rule.

The corpus is not automatically reusable merely because it is available in the
repository. Each item receives an owner, source hash, licence or written
permission, permitted use, redistribution status, retention class, and consent
date. Material marked pending or restricted may support an internal fixture only
under its owner's permission; it cannot enter a published benchmark, external
reader packet, or commercial training set until the rights record is cleared.

The engineering criteria are exact where safety permits exactness. The
protected-source core is the week-six gate; the bounded authoring extension is
built only if that gate passes. The normative command, runtime, security,
schema, budget, rights, and exception rules live in
Appendix~\ref{app:implementation-contract}.

\begin{enumerate}
\item an incomplete critical brief or blueprint blocks model invocation and
      long-form drafting;
\item every declared protected object is either located or explicitly marked
      unparsed; no fixture receives a false protected pass;
\item every seeded change of sign, exponent, unit, number, citation identity,
      label target, quotation, table denominator, or declared claim scope is
      reported at the correct version and source location;
\item capture and inspect leave the source tree byte-identical;
\item a second operator reproduces the canonical PDF hash or receives a
      recorded environment difference that prevents promotion;
\item a preflight critic remains in the release path only when its held-out
      false-negative and false-positive behavior is reported against human
      labels; and
\item a diagnostic remains in the author queue only when its held-out review
      time is lower than the time spent finding and resolving the same true
      issues under the incumbent process.
\end{enumerate}

The seventh criterion replaces phrases such as ``fires heavily'' or ``low enough
not to annoy an editor.'' Let $B_k$ be the minutes spent triaging rule family
$k$ on held-out documents and $U_k$ the incumbent minutes spent locating and
resolving the same adjudicated issues. Retain the family only when
\begin{equation}
  \operatorname{median}(U_k-B_k)>0,
  \label{eq:diagnostic-burden-rule}
\end{equation}
with the complete distribution and the small sample reported. The inequality is
an operating rule, not a population claim. Any silent protected-object failure
overrides it.

\section{The work in calendar order}

Weeks 1--2 freeze the live-case brief, evidence boundary, blueprint schema,
incumbent process, corpus split, protected-object vocabulary, and parser
fixtures. Weeks 3--5 implement planning, bounded drafting, capture, and compare
for LaTeX. Weeks 4--7 overlap the adapter benchmark, mutation harness, and
preflight calibration. Weeks 6--9 connect cause-specific repair, author
dispositions, and the local reader task. Weeks 10--12 run the live feasibility
case, account for every failed gate and minute of burden, and price the
comparative study. Figure
\ref{fig:twelve-week-schedule} and Table
\ref{tab:implementation-sequence} are the single calendar for these intervals.

The evidence review proceeds alongside the build without pretending that
automation completes scholarly judgment. The research lead imports the four
new cornerstone sources and all snowball candidates, while two people screen
load-bearing records independently. The domain reviewer appraises source-to-
claim fit; a different reviewer adjudicates disputed product findings. Package
adapters receive held-out tests before they enter the live path. External
review remains a named dependency, not a status that the implementation team
can award itself.

\section{Deliverables after twelve weeks}

The sponsor receives working software, not a promise of a platform: a
reproducible LaTeX authoring path from brief to release; a blueprint and bounded
writer contract; a locked tuning, mutation, and held-out corpus; package and
critic scorecards with failures and operator burden; one complete feasibility
record; the observed baseline inputs for the economic model; and a proceed,
repair, or stop memorandum. The comparative study is specified and priced, not
smuggled into the feasibility claim.

Continue when the system protects the source, authors can use the findings at
net lower burden, the reader task produces an interpretable outcome, and the
expected value of a comparison exceeds its cost. Repair when one replaceable
adapter or interface fails. Stop the product direction when protected
correspondence cannot be trusted, the reader outcome cannot be observed, or the
workflow consistently adds attention rather than saving it.


\section{Conclusion}
\label{sec:conclusion}

The record, the literature, and the tool field agree on more than we had any
right to expect. Our five projects learned by expensive trial and error that
writing failures arrive in layers and that review machinery cannot write; the
research literature explains why the models write this way and shows that
targeted editing, population-level measurement, and disciplined judging work
where detection and word bans do not; and the open-source field has
independently converged on the same division of labor, with pattern lists for
editors, deterministic counters for pipelines, and honest warnings against
using any of it as an authorship verdict.

What no source establishes, including this one, is that the assembled toolkit
reduces the number of human-feedback rounds a document needs. That is the
only outcome worth the project's budget, it is measurable, and work packages
three and five exist to measure it. The survey's own history sets the
benchmark and the caution in one stroke: its first version passed every
mechanical check we have and failed its only reader. The reader remains the
gate. Everything this proposal builds is machinery for wasting less of the
reader's time.

\bibliography{humanvoice_survey}

\end{document}
